\documentclass[a4paper, oneside]{article}
\usepackage{wrapfig}
\usepackage{graphicx}
\usepackage{amsthm}
\usepackage{amsmath}
\usepackage{amssymb}
\usepackage[a4paper,
            bindingoffset=0.2in,
            left=2cm,
            right=2cm,
            top=2cm,
            bottom=2cm,
            footskip=.25in]{geometry}
\usepackage[italian]{babel}
\usepackage{pgfplots}
\usepackage{tabularx}
\usepackage{tikz-3dplot}
\usepackage{color}
\usepackage{multicol}
\usepackage{arydshln}
\usepackage{mathtools}
\usepackage{enumerate}
\usepackage{graphicx}
\usepackage{svg}
\usepackage{cancel}
\usepackage[d]{esvect}
\usepackage[dvipsnames]{xcolor}
\usepackage{pgfplots}
\usepackage{pifont}
\usepackage{circuitikz}[american voltages]
\usetikzlibrary{patterns}
\makeindex
%\usepackage{animate}
%\usepackage{xfp} % utile se vuoi fare calcoli aggiuntivi
\pgfplotsset{compat=1.18}
\usetikzlibrary{tikzmark}
\newcommand{\TikzNCbar}[4][10pt]{
\tikz[overlay,remember picture]{\draw[#2] (#3) --++(0,-#1) -| (#4);}}
\newcommand{\E}{\mathcal{E}}
\newcommand{\fem}{f_\text{em}}

\graphicspath{ {images/} }

\definecolor{redish}{rgb}{255, 0, 30}
\definecolor{page}{rgb}{0.129,0.157,0.212}
\pagecolor{page}
\color{white}   
\graphicspath{ {./images/} }
\usetikzlibrary{shapes.geometric}
\usetikzlibrary{datavisualization}
\usetikzlibrary{datavisualization.formats.functions}
\pgfplotsset{width=10cm,compat=1.9}

\setlength\dashlinedash{0.2pt}
\setlength\dashlinegap{1.5pt}
\setlength\arrayrulewidth{0.3pt}

\newcommand\eqq{\stackrel{\mathclap{\normalfont\mbox{?}}}{=}}
\newcommand\bulletout  {\labelitemfont \textbullet}
\newcommand{\tab}{\hspace*{2em}}
\newcommand{\xmark}{
\tikz[scale=0.23] {
    \draw[line width=0.7,line cap=round] (0,0) to [bend left=6] (1,1);
    \draw[line width=0.7,line cap=round] (0.2,0.95) to [bend right=3] (0.8,0.05);
}}
\newcommand{\cmark}{
\tikz[scale=0.23] {
    \draw[line width=0.7,line cap=round] (0.25,0) to [bend left=10] (1,1);
    \draw[line width=0.8,line cap=round] (0,0.35) to [bend right=1] (0.23,0);
}}
% Comando:
%   \potato[opzioni]{(x,y)}{scala}
%
% Opzioni = facoltative (es. fill=red!20, draw=black, thick)
% (x,y)   = centro della patata
% scala   = fattore di scala
%
\def\potatoshape{
  (1,0) (2,1.5) (1.6,3) (0.3,2.7) (-0.4,1.2)
}
\newcommand{\potato}[3][draw=white]{
  \begin{scope}[shift={#2}, scale=#3]
    \draw[#1]
      plot [smooth cycle, tension=1]
      coordinates {\potatoshape};
  \end{scope}
}
 \newcommand{\hookbox}[1]{
\begin{center}
\hfill\break
\begin{tikzpicture}
\node[inner sep=0pt,outer sep=0pt,anchor=base] (A) {
\begin{minipage}{\dimexpr\linewidth-5em}
\centering
#1
\end{minipage}
};
% Draw the left bracket
\draw ([xshift=0pt]A.north west) -- ++(0, 0.5) -- ++(0.4, 0);
% Draw the right bracket
\draw ([xshift=0pt]A.south east) -- ++(0, -0.5) -- ++(-0.4, 0);
\end{tikzpicture}
\end{center}} 
\title{Laboratorio di fisica II}
\author{Gariboldi Alessandro}
\date{ }


\begin{document}

\newtheoremstyle{theoremEnv}
                {}          % Space above
                {}          % Space below
                {\slshape}  % Body font
                {}          % Indent amount
                {\bfseries} % Head font
                {.}         % Punctuation after head
                {\newline}         % Space after theorem head
                {}          % Theorem head spec
\theoremstyle{theoremEnv}

\newtheorem{definition}{Definizione}[section]
\newtheorem{theorem}{Teorema}[section]
\newtheorem{lemma}{Lemma}[section]
\newtheorem{observation}{Oss.}[section]
\newtheorem{corollary}{Corollario}[theorem]
\newtheorem{example}{Esempio}[section]
\newtheorem{problem}{Problema}[section]
\newtheorem{solution}{Soluzione}[section]
\newtheorem{proposition}{Proposizione}[section]


\maketitle
\section{23/02/26}
Definiamo le varie quantità che misureremo, e partiamo dalla legge di Ohm:
\begin{gather*}
    V = R I
\end{gather*}
Le misure che facciamo in laboratorio sono macroscopiche anche se si parla di cariche di elettroni che sono quantità microscopiche.\\
La carica dell'elettrone è stata definita esattamente (nel 2019) come:
\begin{gather*}
    e = 1.602176634 \times 10^{-19} C
\end{gather*}
Da quì si definisce l'ampere come:
\begin{gather*}
    1 A = \frac{1 C}{1 s}
\end{gather*}
La corrente la definiamo come:
\begin{gather*}
    I = \frac{dq}{dt}
\end{gather*}
Ovvero quanta carica passa in un certo intervallo di tempo, ma non si guardano i singoli elettroni ma si misura una quantità macroscopica.\\
Questo è vero ma non del tutto, per capire meglio è comodo introdurre il concetto di \textbf{shot noise}: sostanzialmente invece di pensare agli elettroni 
come un flusso omogeneo che può essere associao ad un rumore bianco, si considerano come pacchetti ben definiti 
di carica finita che si può ricondurre se si vuole al rumore della pioggia, da quì shot noise. Questo perchè ogni elettrone ha una carica 
definita ben precisa e siccome passano in momenti leggermente diversi ma casuali, seguono le leggi della statistica e si hanno delle fluttuazioni attorno al valore medio.\\
Quindi l'aspetto microscopico è comunque importante e bisogna ricordarsene.\\
Tutte le tecniche moderne per misurare le correnti sotto lo shot noise, sono le tecniche che vanno al di sotto del singolo elettrone.\\
In generale bisogna sempre avere in mente il limite delle misure che si stanno eseguendo.\\
\subsection{atto pratico}
Supponiamo di avere un conduttore e abbiamo degli ioni (+) che sono posti ad una certa distanza e alcuni elettroni liberi di muoversi, associati ognuno ad uno ione.
\begin{center}
    \begin{tikzpicture}
        \draw[thick] (0,0) sin (0.5,1) cos (1,0) sin (1.5,-1) cos (2,0) sin (2.5,1) cos (3,0) sin (3.5,-1) cos (4,0) sin (4.5,1) cos (5,0);
            \foreach \x in {1.5, 3.5} {
                \node[circle, fill=red!50, minimum size=0.3cm] at (\x,0) {+};
            }
            \foreach \x in {2.5, 4.5} {
                \node[circle, fill=blue!50, minimum size=0.3cm] at (\x,1.2) {-};
            }
            \draw[->, thick] (5,-0.5) -- (5.5,-0.5);
    \end{tikzpicture}
\end{center}
Se prendiamo un cavo di rame con una densità di $d=9 g/cm^3$, e quindi $63,5 g/mol$, avremo una "densità di elettroni su $cm^3$" di:
\begin{gather*}
    n = \frac{N_A \cdot d}{M} = 8.5 \times 10^{22} el/cm^3
\end{gather*}
Questo vuol dire che in un $\mu m$ cubo di rame ci sono circa $10^{11}$ elettroni liberi.\\
Ora si capisce bene il concetto di densità di carica che è:
\begin{gather*}
    \rho = n \cdot e
\end{gather*}
Ovvero la quantità di carica per unità di volume, poichè si ricorda che $e$ è la carica di un elettrone e $n$ il numero di elettroni liberi per unità di volume.\\
Questo non è sempre vero perchè non tutti gli elettroni sono liberi, ma è una buona approssimazione.\\
\subsection{cos'è una carica}
Considero una carica $q$ in un campo elettrico $\vv{E}$, questa carica subisce una forza data da:
\begin{gather*}
    \vv{F} = q \vv{E}
\end{gather*}
Se ho un campo elettrico con una carica all'interno, questo campo applicherà una forza alla mia carica e quest'ultima subirà una certa accelerazione, quindi si potrebbe naturnalmente pensare che la velocità della carica dipenda anche dal tempo.\\
Questo dilemma viene risolto dalla legge di Ohm, che mi dice che la corrente è proporzionale alla tensione, e quindi se la tensione è costante allora la corrente è costante, e quindi la carica che passa in un certo intervallo di tempo è costante.\\
Infatti gli elettroni non sono completamente liberi di muoversi ma sbattendo tra loro subiscono una certa frizione, meglio detta \textbf{resistenza} così facendo la velocità si stabilizza dopo una certa quantità di tempo, e quindi si ha una velocità media che è costante.\\
\begin{example}
    Prendiamo un gas a una temperatura $T$ costituito di elettroni, avremo che $\frac{1}{2} m v_t^2 = \frac{3}{2} k_B T$ è una relazione che ci dice che la velocità termica è proporzionale alla temperatura, e quindi se la temperatura è alta allora la velocità termica è alta.\\
    Questa velocità termica mi dice la velocità a cui gli elettroni si muovono mediamente ad una certa temperatura, costantemente accelerando e decelerando urtando tra di loro (è comodo immaginarlo come un elettrone che vibra si muove molto velocemente ma non si sposta nel materiale più di tanto).\\
    Calcolando questa velocità termica per una temperatura di $300 K$ (temp. ambiente) ottengo:
    \begin{gather*}
        v_t = \sqrt{\frac{3 k_B T}{m}} = \sqrt{\frac{3 k_b T}{n_e}} = \sqrt{\frac{\cancel{3} 1,3810^{-23} T/k \cancel{3}00 k}{\cancel{9},1 10 ^{-31}}kg} = 1,16 \times 10^5 m/s
    \end{gather*}
    Ricordiamo che il moto degli elettroni in un metallo è quantistico ma classicamente possiamo dire che da un moto di agitazione termica di velocità di centinaia di km/s (la stessa velocità a cui viaggiano gli atomi nell'aria), quindi la corrente viaggia più o meno alla stessa velocità?\\ 
\end{example}

\subsection{Velocità di deriva}
Consideriamo gli elettroni all'equilibrio termodinamico e quindi con velocità termica $v_t$ ma moto medio nullo (stanno vibrando ma non si stanno spostando).\\
Ora ci applichiamo un campo elettrico e voglio trovare una relazione che mi descriva come si muovono gli elettroni in un materiale, generalmente in un conduttore.\\
Consideriamo ora che gli elettroni che si muovono in un corpo, considero i loro urti elastici in un intervallo generale $\Delta t$ e in questo intervallo avrò una certa velocità $\Delta v$ sto quindi considerando un'accelerazione.\\
In sintesi ho un elettrone fermo che subisce una forza da parte del campo elettrico, e quindi subisce un'accelerazione, e dopo un certo intervallo di tempo $\Delta t$ subisce un urto elastico e quindi si ferma e ne fa partire un'altro, questo processo continua a ripetersi continuamente.\\
\begin{gather*}
    \Delta \vec{v} = \vec{a} \Delta t = \frac{\vec{F}}{m} \Delta t = \frac{-e \vec{E}}{m} \Delta t\\
    \Delta \vec{v} = -\frac{e \vec{E}}{m} \Delta t
\end{gather*}
Con $\Delta t = \frac{\Delta l}{v_t}$ Con $\Delta l$ che è il cammino medio in quell'intervallo di tempo, con questo possiamo calcolare la velocità di deriva $v_d$ che è la velocità media degli elettroni in un conduttore, e quindi:
\begin{gather*}
    \vec{v_d} =-\frac{e \vec{E}}{2m} \Delta t
\end{gather*}
Ma $\Delta t= \frac{\Delta l}{v_t}$ e siccome $v_d \ll v_t$
\begin{gather*}
    \vec{v_d} = - \frac{e \vec{E}}{2m} \frac{\Delta l}{v_t} = k(t) \vec{E}
\end{gather*}
Ricordiamo che $v_d := \left\langle \Delta v\right\rangle_T$ è la velocità media degli elettroni liberi in un conduttore.
verifichiamo ora l'ipotesi di $v_d \ll v_t$
\subsection{densità di corrente}
Si parlerà dell'equazione di continuità che deriva dal principio di conservazione della carica.\\
E si arriverà alla prima legge di kirkhoff.\\
\hfil\\
Prendiamo un elettrone in concizioni stazionarie (che non vuol dire che la corrente è stazionaria che vuol dire che $\vec{E} = 0$ cosa che comunque considereremo per parte di questo corso).\\
Ora definiamo la densità di corrente $\vec{J}$ come la quantità di carica che passa attraverso una superficie $S$ per unità di tempo (il flusso di un campo su una superficie), e quindi la corrente è data da:
\begin{gather*}
    I = \int_S \vec{J} \cdot \vec{dS} = \int_S \vec{J} \widehat{n} dS
\end{gather*}
Dove $\vv{J}$ è proprio la densità di corrente, e $\widehat{n}$ è il versore normale alla superficie $S$.\\
Possiamo anche osservare che la corrente è data da $I = n q v_d S$, e quindi possiamo scrivere la densità di corrente come:
\begin{gather*}
    \vec{J} = n q \cdot \vec{v_d} \propto \vec{E}
\end{gather*}
Da questo si vede che vale sempre che $\vec{J} \parallel \vec{E}$ indipendentemente dalla carica, per convenzione si dice che la corrente si "muove" dal polo positivo 
a quello negativo ma di fatto sono gli elettroni che agitandosi generano la corrente.\\
$\vec{J}$ è una densità di corrente, e si misura in $A/m^2$, come mostrato dall'equazione dimensionale seguente.\\
\begin{gather*}
    [J] = \frac{C}{m^3} \frac{m}{s} = \frac{A}{m^2} 
\end{gather*}
\section{25/02/26}
La scorsa lezione abbiamo detto che possiamo avere delle cariche $n \sim 9 \times 10^{22} \text{ atomi/cm}^3$, che possono avere una certa densità di carica $\rho$\\
e che possono essere "accelerate" da un campo elettrico $\vv{E}$ ma che sotto le nostre ipotesi di $v_d \ll v_t$ e di urti che avvengono in un intervallo di tempo $\Delta t$ possiamo dire che $\vv{v_d} \propto \vv{E}$.\\
Inoltre abbiamo visto che la corrente è data da $I = \int_S \vv{J} \cdot \hat{n} dS$ e che $\vv{J} = n q \vv{v_d} \parallel \vv{E}$, quindi possiamo dire che $I \propto E$.\\
Come ultima cosa siamo arrivati alla legge di Ohm in forma locale:
\begin{gather*}
    \vv{J}(\vv{r}) = \sigma_c \vv{E}(\vv{r})
\end{gather*}
Dove $\sigma_c$ è la conducibilità elettrica, che è una proprietà del materiale.\\
L'unica cosa che avevamo dimostrato era la velocità di deriva, ma avevamo visto che $\vv{v_t}$ è molto grande circa $100 km/s$\\
Altro aspetto importante è che la legge di ohm si può riscrivere come:
\begin{gather*}
    \vv{E} = \rho_c \vv{J}
\end{gather*}
Dove $\rho_c$ è la resistività elettrica, che è l'inverso della conducibilità ovvero $\rho_c = \frac{1}{\sigma_c}$.\\
\begin{example}[Calcolo della velocità di deriva]
    Ponendo di avere una corrente di un ampere che passa attraverso un filo di raggio
    un millimetro. Dunque la sezione $S = \pi r^{2} = \pi$ mm$^{2}$. Dunque la corrente è:
    \begin{gather*}
        I = 1 \ A = n q \cdot v_d \cdot S  \ \Longrightarrow \ v_d = 2.5 \cdot 10^{-5} \frac{m}{s}
    \end{gather*} 
\end{example}
\subsection{Dal potenziale elettrico alla prima legge di Kirchhoff}
In fisica la carica si conserva sempre, noi abbiamo definito la corrente la carica su unità di tempo $I = \frac{dQ}{dt}$.\\
Ora prendiamo una carica (uscente quindi con segno negativo) $Q$ che attraversa una superficie chiusa $S$.
Se considero la carica che esce come $-Q(t)$ in base al tempo quindi in un tempo $dt$ avrò una diminuzione della carica:
\begin{gather*}
    -dQ(t) = \int_S \vv{J} \cdot \vv{dS} dt
\end{gather*}
\begin{gather*}
    I = \frac{dQ(t)}{dt} = \int_S \vv{J} \cdot \vv{dS}
\end{gather*}
Considero che la corrente è uscente, quindi $I = -\frac{dQ(t)}{dt}$ e quindi:
\begin{gather*}
    -\frac{dQ(t)}{dt} = \int_S \vv{J} \cdot \vv{dS}
\end{gather*}
Ora si applica il teorema della divergenza che ci dice che l'integrale di un campo vettoriale su una superficie chiusa è uguale all'integrale della divergenza di quel campo vettoriale su tutto il volume racchiuso da quella superficie, e quindi:  
\begin{gather*}
    -\frac{dQ(t)}{dt} = \int_S \vv{J} \cdot \vv{dS}\\
    -\frac{dQ(t)}{dt} = \int_V \vv{\nabla} \cdot \vv{J} dV\\
\end{gather*}
Se calcoliamo la divergenza abbiamo che $\vv{\nabla} \cdot \vv{J} = \frac{\partial J_x}{\partial x} + \frac{\partial J_y}{\partial y} + \frac{\partial J_z}{\partial z}$, ricordando però che $\vv{J} = n q \vv{v_d}$, quindi se se consideriamo $n$ e $q$ costanti e che $\vv{v_d}$ è costante nel tempo, allora la divergenza è nulla, poichè stiamo considerando la derivata rispetto al tempo di ogni termine che è costante nel tempo, e quindi:
\begin{gather*}
    \vv{J} = n q \vv{v_d} \Rightarrow \vv{\nabla} \cdot \vv{J} = n q \vv{\nabla} \cdot \vv{v_d}\\
    \vv{\nabla} \cdot \vv{J} = nq\frac{\partial v_d}{\partial x} + nq\frac{\partial v_d}{\partial y} + nq\frac{\partial v_d}{\partial z} = 0
\end{gather*}
Ma $\vv{v_d} = \begin{pmatrix}
    \frac{d r_x}{dt}\\
    \frac{d r_y}{dt}\\
    \frac{d r_z}{dt}
\end{pmatrix}$ e quindi se considero che $\frac{d r_x}{dt}$ è costante nel tempo, allora $\frac{\partial v_d}{\partial x} = 0$ e così per gli altri due termini, e quindi:
\begin{gather*}
    -\frac{dQ(t)}{dt} = \int_V \frac{\partial \rho(x,y,z,t)}{\partial t} dx dy dz\\
    -\int_V \frac{\partial \rho}{\partial t} dV = \int_S \vv{J} \cdot \vv{dS} = \int_V \vv{\nabla} \cdot \vv{J} dV\\
    -\frac{\partial \rho}{\partial t} = \vv{\nabla} \cdot \vv{J} \Rightarrow \boxed{\vv{\nabla} \cdot \vv{J} + \cancel{\frac{\partial \rho}{\partial t}} = 0}\\
    \Rightarrow \boxed{\vv{\nabla} \cdot \vv{J} = 0} \Rightarrow \boxed{\int_{S_\text{chiusa}} \vv{J} \cdot \vv{dS} = 0}
\end{gather*}
\subsection{nodi}
Un nodo è un punto di incontro tra due o più conduttori.
La somma algebrica delle correnti che entrano in un nodo è uguale alla somma algebrica delle correnti che escono da un nodo.\\
Se ho una superficie chiusa che racchiude il nodo $n$ considerando $\int_S \vv{J} \widehat{n} dS = 0$
\begin{center}
    \begin{tikzpicture}
        \draw[dashed] (0,0) circle (1.4);
        %curve intersecanti
        \draw (-1,-1) to [bend left=20] (1,1);
        \draw (-1,1) to [bend right=20] (1,-1);
        \draw(-1,-1) circle (3pt) node[below] {$S_3$};
        \draw(1,1) circle (3pt) node[above] {$S_1$};
        \draw(-1,1) circle (3pt) node[above] {$S_4$};
        \draw(1,-1) circle (3pt) node[below] {$S_2$};
        %versori n uscenti
        \draw[->, thick] (-1,-1) -- (-1.4,-1.4) node[below] {$\hat{n} \quad I_4$}; 
        \draw[->, thick] (1,1) -- (1.4,1.4) node[above] {$\hat{n} \quad I_2$}; 
        \draw[->, thick] (-1,1) -- (-1.4,1.4) node[above] {$\hat{n} \quad I_1$}; 
        \draw[->, thick] (1,-1) -- (1.4,-1.4) node[below] {$\hat{n} \quad I_3$}; 
    \end{tikzpicture}
\end{center}
\begin{gather*}
    \sum_{i=1}^n I_i = 0\\
    I_1 + I_2 + I_3 + I_4 = 0
\end{gather*}
\subsection{Maglie e rami}
\begin{center}
    \begin{tikzpicture}[scale = 0.5]
        \draw (-1,-1) to [bend left=20] (1,1);
        \draw (-1,1) to [bend right=20] (1,-1);\
        \filldraw(-0.4,0) circle (3pt) node[below] {$N$};
    \end{tikzpicture}
    $\qquad\qquad$
    \begin{tikzpicture}
        \draw (-1,0) -- (1,0) node[midway, above] {$R$};
        \filldraw (-1,0) circle (2pt);
        \filldraw (1,0) circle (2pt);
        \draw (-1,0) -- (-1.5,-1);
        \draw (-1,0) -- (-1.5,1);
        \draw (1,0) -- (1.5,-1);
        \draw (1,0) -- (1.5,1);
    \end{tikzpicture}
    $\qquad\qquad$
    \begin{tikzpicture}
        \draw (0,0) rectangle (2,2) node[midway] {$M$};
        \filldraw (0,0) circle (2pt);
        \filldraw (2,0) circle (2pt);
        \filldraw (0,2) circle (2pt);
        \filldraw (0,2) circle (2pt);
        \filldraw (2,2) circle (2pt);
    \end{tikzpicture}
\end{center}
Un ramo è un conduttore che collega due nodi, una maglia è un percorso chiuso che non contiene altri nodi al suo interno.\\
Più maglie formano una rete.\\
\subsection{Il potenziale elettrostatico}
Se in un campo il rotore è uguale a 0 vuol dire che è conservativo, nel nostro caso il rotore è proprio:
\begin{gather*}
    \vv{\nabla} \land \vv{E} = 0
\end{gather*}
\begin{wrapfigure}{r}{0.4\textwidth}
    \begin{center}
    \begin{tikzpicture}
        %assi
        \draw[->] (0,0) -- (3,0) node[below] {$x$};
        \draw[->] (0,0) -- (0,3) node[left] {$y$};
        \draw[->] (0,0) -- (-1,-1) node[below] {$z$};
        \draw[->] (0,0) -- (2,1) node[midway, above] {$\vv{r_A}$};
        \filldraw (2,1) circle (1pt) node[above] {$A$};
        \draw[->] (0,0) -- (1,-0.5) node[midway, below] {$\vv{r_B}$};
        \filldraw (1,-0.5) circle (1pt) node[below] {$B$};
        \draw[->, cyan] (2,1) to [bend left=20] (1,-0.5);
        \node[cyan] at(2.1,0.5) {$\Gamma$};
        \node at(0,0) [left] {$O$};
    \end{tikzpicture}
\end{center}
\end{wrapfigure}
Consideriamo una carica $Q$ posta in $O$, e vogliamo calcolare quanto lavoro fa il campo elettrico generato da $Q$ su una carica di prova $q$ che si muove da $A$ a $B$ lungo un percorso $\Gamma$.\\
Il campo elettrico generato da $Q$ sarà dato da:
\begin{gather*}
    \vv{E}(\vv{r}) = \frac{Q}{4 \pi \varepsilon_0 r^2} \hat{r}
\end{gather*}
Il lavoro che fa il campo elettrico su $q$ è dato da:
\begin{gather*}
    \mathcal{L}_{AB}^\Gamma = \int_\Gamma \vv{F} \vv{dl} = \int_\Gamma q \vv{E} \vv{dl} = \frac{qQ}{4 \pi \varepsilon_0} \int_{r_a}^{r_b}  \frac{1}{r^2} \overbrace{\widehat{r} \widehat{r}}^{=1} dr =\\
    \int_{r_a}^{r_b} \frac{1}{r^2} dr = -\frac{1}{r}\Big|_{r_a}^{r_b} = -\frac{1}{r_b} + \frac{1}{r_a}
\end{gather*}
\begin{gather*}
    \mathcal{L}_{AB}^\Gamma = \frac{qQ}{4 \pi \varepsilon_0} \left( -\frac{1}{r_b} + \frac{1}{r_a} \right)
\end{gather*}
Quest'ultimo è il principio di conservatività del campo (1a legge di Kirchhoff).\\
Cerchiamo ora di definire il potenziale elettrico. \\
Per farlo possiamo considerare il lavoro che fa il campo elettrico sulla carica $q$, questa carica è quella di cui calcoliamo il lavoro generato da $Q$, ma siccome a sua volta ha una carica 
che potrebbe influenzare il campo elettrico generato da $Q$, dobbiamo normalizzare $U(r)$ per $q$ ($\frac{U}{q}$) e consideriamo $q$ come infinitesima così da non influenzare 
il campo elettrico circostante. Così facendo possiamo misurare quanto lavoro farebbe il campo in un determinato punto:
\begin{gather*}
    U(r) = \frac{q \ Q}{4 \pi \varepsilon_0} \frac{1}{r} \Rightarrow \lim_{q \to 0} \frac{U}{q} = \boxed{V(r) = \frac{Q}{4 \pi \varepsilon_0} \frac{1}{r} + cost.}\\
\end{gather*}
Abbiamo ottenuto il potenziale del campo, di fatto su una carica puntiforme, in un punto $P$ ben definito a una distanza $r$ dalla sorgente $Q$.$V(r)$ rappresenta così quanto lavoro farebbe una carica in ogni punto, se ad una distanza $r$ il potenziale vale $10 V$ allora una carica di $1 C$ farebbe un lavoro di $10 J$.\\
Da notare che quì abbiamo riesplcitato la costante di integrazione, nei casi precedenti la avevamo omessa perchè rappresenta il fatto che il potenziale è definito a meno di una costante arbitraria, di solito si usa come riferimento della carica di prova l'infinito, e quindi si pone $V(\infty) = 0$ e quindi $cost. = 0$. 
Ma quando abbiamo fatto il limite e dividiamo per $q$ che tende a 0 dobbiamo esplicitarla poichè altrimenti $U$ schizzerebbe all'infinito.\\
Possiamo ora calcolare il lavoro tra due punti $A$ e $B$ del campo $\vv{E}$ lungo un percorso generico e calcolare quella che si chiama differenza di potenziale:
\begin{gather*}
    V_A - V_B = \int_{A}^{B} \vv{E} \cdot \vv{dl}
\end{gather*}
Con la differenza di potenziale che è proprio $V_A - V_B $ e si misura in volt $\frac{J}{C}$ ovvero lavoro su unità di carica.\\
\subsection{Seconda legge di Kirchhoff}
Ricordiamo che se abbiamo una curva chiusa $\Gamma$ allora:
\begin{gather*}
    \oint_\Gamma \vv{E} \cdot \vv{dl} = 0
\end{gather*}
\begin{center}
    \begin{tikzpicture}
        \draw (-1,-1) rectangle (1,1);
        \filldraw (1,1) circle (1pt) node[above] {$C$};
        \filldraw (-1,-1) circle (1pt) node[below] {$A$};
        \filldraw (-1,1) circle (1pt) node[above] {$B$};
        \filldraw (1,-1) circle (1pt) node[below] {$D$};
        \draw[dashed] (0,0) circle (1.4);
        \node at(-1,0) [right] {$R_1$};
        \node at(1,0) [left] {$R_3$};
        \node at(0,1) [below] {$R_2$};
        \node at(0,-1) [above] {$R_4$};
        \node at(0,-1.4)[below] {$\Gamma$}; 
    \end{tikzpicture}
\end{center}
Poichè se calcoliamo il lavoro su una curva chiusa otteniamo 0, allora se consideriamo la curva chiusa $\Gamma$ che passa per i punti $A$, $B$, $C$ e $D$, avremo che:
\begin{gather*}
    \oint_\Gamma \vv{E} \cdot \vv{dl} = \int_{A}^{B} \vv{E} \cdot \vv{dl} + \int_{B}^{C} \vv{E} \cdot \vv{dl} + \int_{C}^{D} \vv{E} \cdot \vv{dl} + \int_{D}^{A} \vv{E} \cdot \vv{dl} = 0\\
    (V_A - V_B) + (V_B - V_C) + (V_C - V_D) + (V_D - V_A) = 0
\end{gather*}
\subsection{parliamo un po' di più della resistività}
Partendo da $\vv{E} = \rho_R \vv{J}$ e considerando l'unità di misura della resistività:
\begin{gather*}
    [\rho_R] = \left[\frac{N}{C} \frac{m^2}{A}\right] \\
    \left[ \frac{J \ m}{C \ A} \right] = \left[ V \frac{m}{A} \right] = \left[ \Omega m\right]   
\end{gather*}
Prendiamo l'esempio del traliccio dell'alta tensione, dove alcune mucche morivano, questo succede perchè anche se ovviamente il cavo è la parte che fa meno resistenza e il traliccio ha un punto di materiale isolante, comunque una parte 
della corrente passa attraverso il traliccio e quindi una piccola parte della corrente passa attraverso il traliccio, questo accadeva con le mucche per via dei passi molto lunghi, quindi 
facevano si di avere una differenza di potenziale tra le due zampe, e quindi una corrente che passava attraverso il loro corpo, e questo era sufficiente per ucciderle.\\
\hfill\\
Altra cosa è il fatto che la differenza di potenziale tra il terreno e le nuvole e circa $200 V$ per $m$, noi non prendiamo la scossa perchè fungiamo da conduttori, ma quando si accumula sulle nuvole comunque da qualche parte devono 
scaricarsi, e quindi quando si accumula una certa quantità di carica sulle nuvole, questa si scarica attraverso un fulmine\\
\subsection{La resistenza elettrica}
\begin{gather*}
    \vv{E} = \rho_R \vv{J}
\end{gather*}
\begin{center}
    \begin{tikzpicture}
        %cilindro
        \draw (-1,0) -- (1,0);
        \draw (-1,1) -- (1,1);
        \draw[pattern = north east lines] (-1,0.5) ellipse (0.3 and 0.5) node[above] {$S$};
        \draw[pattern = north east lines] (1,0.5) ellipse (0.3 and 0.5) node[above] {$S$};
        \node at(0,1) [above] {$\ell$};
        \draw[thin] (-1,0.5) -- (1,0.5) node[at start, below]{\small{$A$}} node[at end, below]{\small{$B$}};
        \filldraw (-1,0.5) circle (1pt);
        \filldraw (1,0.5) circle (1pt);
        \draw[->] (1.1,0.5) -- (1.5,0.5) node[above] {$\widehat{n}$};
    \end{tikzpicture}
\end{center}
Considero un conduttore di sezione $S$ (che consideriamo relativamente piccola a $\ell$) e di lunghezza $\ell$, con una resistività elettrica $\rho_R$ e con una corrente $I$ che lo attraversa, e voglio calcolare la differenza di potenziale tra i due capi del conduttore, ovvero $V_A - V_B$.\\
Consideriamo che localmente $\vv{E} = \rho_R \vv{J}$ e siccome  $S$ è molto piccola rispetto a $\ell$ allora possiamo considerare $E$ e $J$ costanti lungo il conduttore, svolgendo i conti:
\begin{gather*}
    \int_{A}^{B} \vv{E} \cdot \vv{dl} = V_A - V_B = \int_{A}^{B} \rho_R \vv{J} \cdot \vv{dl} = \int_{A}^{B} \rho_R \frac{\vv{J} S}{S} \cdot \widehat{n} dl \qquad \text{qui $\vv{J}$ e $S$ si considerano $\simeq$ cost.}\\
    \Rightarrow I \underbrace{\int_{A}^{B} \frac{\rho_R}{S} dl}_{R} = I \frac{\rho_R \ell}{S} = I R\\
    \Delta V = I R
\end{gather*}
Quindi $R = \frac{\rho_R \ell}{S}$ è la resistenza elettrica, che è una proprietà del conduttore, e si misura in $\Omega$.\\
Dalla legge di ohm $R = \frac{V}{I}$ possiamo pensare alla resistenza come alla differenza di potenziale su unità di corrente, e quindi se $I$ è la quantità di carica che passa in un certo intervallo di tempo, allora $R$ è la differenza di potenziale su unità di carica che passa in un certo intervallo di tempo.\\
In alternativa possiamo pensare che siccome $V = \frac{J}{C}$ ovvero lavoro su unità di carica, allora $R$ è il lavoro su unità di carica su unità di corrente, quindi $R$ è il lavoro su unità di carica su unità di carica su unità di tempo, quindi $R$ è il lavoro su unità di carica al secondo, e quindi $R$ è la potenza per unità di carica.\\
\subsection{applicazione delle resistenze}
Abbiamo in primo luogo che la differenza di potenziale $\Delta V$ è presente solo se, ovviamente c'è corrente, e se non c'è resistenza nulla ($R \neq 0$). Quindi abbiamo 
che se $R = 0$ allora $\Delta V = 0$ e quindi non c'è differenza di potenziale, e se $I = 0$ allora $\Delta V = 0$ e quindi non c'è differenza di potenziale.\\
\begin{center}
    \begin{tikzpicture}
        \draw(0,0) to[short,R, l=$R$] (2,0);
        \filldraw(0,0) circle(1pt) node[below] {\small{$A$}};
        \filldraw(0.4,0) circle(1pt) node[below] {\small{$B$}};
        \filldraw(1.6,0) circle(1pt) node[above] {\small{$C$}};
        \filldraw(2,0) circle(1pt) node[above] {\small{$D$}};
    \end{tikzpicture}
\end{center}
Ad esempio in questa illustrazione abbiamo che tra i punti $B$ e $C$ c'è una resistenza $R$, quindi se tra $B$ e $C$  c'è corrente allora  c'è differenza di potenziale, ma tra i punti $A$ e $B$ e tra i punti $C$ e $D$ non c'è resistenza, quindi se tra questi punti non c'è differenza di potenziale indipendentemente dal passaggio di corrente.\\
Per calcolare la resistenza di un conduttore possiamo usare la formula $R = \frac{\rho_R \ell}{S}$, e quindi se vogliamo calcolare la resistenza di un filo di rame, dobbiamo conoscere la resistività del rame $\rho_{Cu}$ che è circa $1.7 \times 10^{-8} \Omega m$, e poi dobbiamo conoscere la lunghezza del filo $\ell$ prendiamo ad esempio $1 m$ e la sezione del filo $S$ prendiamo ad es $10^{-3} m^2$ ovvero $1mm$.
\begin{gather*}
    R = \frac{1,7 \times 10^{-8} \Omega \cancel{m} 1 \cancel{m}}{\pi ( 10^{-3} \cancel{m})^2}
\end{gather*}

\section{2/3/26}
\subsection{Generatori di forza elettro motrice}
Un generatore di potenziale è un dispositivo che fornisce una differenza di potenziale costante tra i suoi terminali, indipendentemente dalla corrente che lo attraversa.\\
Cerchiamo di capire come si misura la forza elettromotrice (f.e.m.) di un generatore di potenziale.\\
Il generatore ideale viene rappresentato in un circuito in questo modo:
\begin{center}
    \begin{tikzpicture}
        \draw (0,0) to [battery1, invert, l=$\pm$,] (0,2);
            \draw (0,2) to [short] (2,2);
            \draw (2,2) to [short] (2,0);
            \draw (2,0) to [short] (0,0);
            \draw[->](1.5,1.8) -- (1.5,0.2) node[right, midway] {$i$};
            \draw[->](2.5,1.8) -- (2.5,0.2) node[right, midway] {$\vv{E}_s$};
    \end{tikzpicture}
\end{center}
    Per ora non studiamo nel dettaglio come funziona ci limitiamo a considerare che tra il $+$ e il $-$ del generatore c'è una forza che separa le cariche e che va contro il campo elettrostatico $\vv{E}_s$, e questa forza è rappresentata proprio da un altro campo elettrico che chiameremo $\vv{E}_e$, che scopriremo essere proprio la forza elettromotrice, se consideriamo solo il campo elettrostatico sulla curva chiusa $\Gamma$ che rappresenta il nostro circuito, allora avremo che:
    \begin{gather*}
        \oint_\Gamma \vv{E}_s \cdot \vv{dl} =  0
    \end{gather*}
    Poichè questo campo è conservativo. Ora consideriamo la somma dei campi $\vv{E}_e$ ed $\vv{E}_s$, 
    \par
    \begin{wrapfigure}{r}{0.4\textwidth}
            \begin{tikzpicture}
                \draw (0,0) to [generic, invert, l=$\pm$,] (0,1);
                \draw (0,0) to[short] (0,-0.5);
                \draw (0,1) to[short] (0,1.5);
                \draw (0,-0.5) to[short] (2,-0.5);
                \draw (0,1.5) to[short] (2,1.5);
                \draw (2,1.5) to[short, f=$i$] (2,-0.5);
                \draw[<-] (-1,0) -- (-1,1) node[left, at start] {$\vv{E}_e$};
                \node at(1,1) [above] {$\Gamma$};
                \filldraw(0,-0.1) circle (1pt) node[left] {$B$};
                \filldraw(0,1.1) circle (1pt) node[left] {$A$};
                \draw[->] (0.1,0.2) -- (0.1,0.8) node[left, midway] {$i$};
                \draw[->] (0.5,0.1) -- (0.5,1) node[right, midway] {$\vv{E}_s$};
            \end{tikzpicture}
    \end{wrapfigure}\noindent
    Ci interesserà più avanti scoprire queste forze interne al generatore, per il momento però ricordiamo solo che esistono.\\
    La curva chiusa $\Gamma$ rappresenta il mio circuito, ma se consideriamo solo il pezzo che attraversa il generatore di potenziale, allora avremo:
    \begin{gather*}
        \oint_\Gamma \vv{E} \cdot \vv{dl} = \\
        \int_{A}^{B} \vv{E}_e \cdot \vv{dl} + \int_{B}^{A} \vv{E}_s \cdot \vv{dl} = 0
    \end{gather*}
    Questo perchè il campo $\vv{E}_e$ è conservativo se considerato all'interno del generatore, di fatto abbiamo dei vettori dritti che spostano le cariche da - a +, e non abbiamo rotazioni interne, stessa cosa ovviamente per $\vv{E}_s$.\\
    Di fatto se consideriamo solo il pezzo che attraversa il generatore, il pezzo del campo $\vv{E}_s$ all'interno lo chiameremo spesso $\vv{E}_s'$ però si ricorda essere di fatto lo stesso campo.\\
    Se consideriamo però tutta la curva $\Gamma$ abbiamo che $\vv{E}_e$ non è più presente fuori dal generatore quindi ho che compiuta una rotazione sulla curva ho dell'energia che non mi è stata restituita. So però che il campo $\vv{E}_s$ è conservativo anche su 
    tutta la curva e quindi posso considerarlo nullo quando vedo la somma dei due, così facendo posso trovare che quella che è la forza elettromotrice del generatore è quella generata dal campo $\vv{E}_e$ che è in contrapposizione con il campo $\vv{E}_s$ e quindi:
    \begin{gather*}
        \underbrace{\oint_\Gamma \vv{E}_s \cdot \vv{dl}}_{= 0 \text{\small{ perchè conservativo}}} + \vv{E}_e \cdot \vv{dl} = \boxed{\int_{A}^{B} \vv{E}_e \cdot \vv{dl} \neq 0} := \text{f.e.m.}
    \end{gather*}
    Questa è la definizione operativa di forza elettromotrice, ovvero la differenza di potenziale che si misura ai capi del generatore quando non c'è corrente che lo attraversa, altrimenti infulenzerebbe il campo all'interno del generatore.\\
    Di fatto questo è un campo non conservativo, perchè stiamo considerando si un campo di vettori "dritti" ma su una curva più grande di dove lui è definito facendo si che il rotore sia non nullo.\\
    \subsection{Generatore reale schematizzato}
    Di fatto un generatore nella realtà non riesce a separare le cariche senza disperdere nessuna energia, quindi in realtà un generatore di potenziale è sempre accompagnato da una resistenza interna, che è la resistenza che si oppone al passaggio della corrente all'interno del generatore stesso.\\
    Possiamo quindi schematizzare un generatore di potenziale reale in questo modo:
    \begin{center}
        \begin{tikzpicture}
            \draw (0,0) to [battery1, invert, l=$\mathcal{E}_x f_x \quad \pm$,] (0,2);
            \draw (0,2) to [short] (2,2);
            \draw (2,2) to [R, l=$r_x \quad \rho_x$] (2,0);
            \filldraw(2,0) circle (2pt) node [below] {$A$};
            \filldraw(0,0) circle (2pt) node [below] {$B$};
        \end{tikzpicture}
    \end{center}
    Con $\mathcal{E}_x$ che rappresenta la forza elettromotrice del generatore, $f_x$ che è il modulo del campo che precedentemente abbiamo chiamato $\vv{E}_s$ e rappresenta la forza del campo per unità di carica quindi se si vuole l'efficienza del generatore, $r_x$ che rappresenta la resistenza interna del generatore e $\rho_x$ che rappresenta la resistività elettrica associata alla resistenza.\\
    Schematizzando così il nostro generatore abbiamo in questo caso che $i=0$ e quindi possiamo definire la f.e.m. come proprio la differenza di potenziale ai capi del circuito senza passaggio di corrente.
    \begin{gather*}
        f.e.m. = V_A - V_B
    \end{gather*}
    Ovviamente a noi del passaggio di corrente serve per misurare questa grandezza, quindi come possiamo fare?\\
    Per misurare la f.e.m. sperimentalmente sfrutteremo l'andamento della corrente, fino ad un certo punto è circa lineare 
    ma da un certo punto in poi diventa non lineare, questo punto in cui si ha un salto alla non linearità, può essere sfruttato per misurare la f.e.m. del generatore.\\
    Possiamo visualizzare questo andamento in questo modo:
    \begin{center}
        \begin{tikzpicture}
            %assi x y
            \draw[->] (0,0) -- (3,0) node[below] {$i$};
            \draw[->] (0,0) -- (0,3) node[left] {$V_A-V_B$};
            %linea
            \draw[thick] (0,2) -- (2,0.5) node[midway, above] {$r$} node[at start, left]{$f$};
            \draw[thick, cyan] (2,0.5) to[bend left = 10] (2.5,0);
            \node at(2.5,0) [below] {$i_m$};
        \end{tikzpicture}
    \end{center}
    Sfruttando la relazione precedente e integrando ciò che abbiamo detto finora si puo dire:
    \begin{gather*}
        \Delta V = V_A - V_B = f_\text{em} - r i
    \end{gather*}
    Ovvero la differenza di potenziale che si misura ai capi del generatore è data dalla forza elettromotrice meno la caduta di potenziale dovuta alla resistenza interna.\\
    Questo torna con la situazione precedente perchè quando la corrente è molto piccola, la caduta di potenziale dovuta alla resistenza interna è trascurabile, e quindi $V_A - V_B \simeq f_\text{em}$, ma quando la corrente aumenta, la caduta di potenziale dovuta alla resistenza interna diventa significativa, e quindi $V_A - V_B$ diminuisce.\\
    Questo ci dice anche che all'aumentare della corrente l'effetto della resistenza sale fino ad un valore massimo che è dove la differenza di potenziale non è più in grado di tenerlo, questo è praticamente il limite del generatore ed è dove comincia a surriscaldarsi, e proprio per questo si ha il salto alla non linearità.\\
    \par
    \begin{wrapfigure}{r}{0.4\textwidth}
    \begin{center}
    \begin{tikzpicture}
            %assi x y
            \draw[->] (0,0) -- (3,0) node[below] {$i$};
            \draw[->] (0,0) -- (0,3) node[left] {$V_A-V_B$};
            %linea
            \draw[thick] (0,2) -- (2,0.5) node[midway, above] {$r$} node[at start, left]{$f$};
            \draw[thick, cyan] (2,0.5) to[bend left = 10] (2.5,0);
            \node at(2.5,0) [below] {$i_m$};
            \draw[dashed] (0,2.5) -- (2.5,2.5);
            \draw[dashed] (2.5,0) -- (2.5,2.5); 
            \node[orange] at(2.5,2.5)[above right] {$W$};
            \draw[orange] (0,2.5) ..controls(2.5,2.5).. (2.5,0);
        \end{tikzpicture}
    \end{center}
    \end{wrapfigure}\noindent
    Posso visualizzarla anche in un'altro modo.\\
    Considero la potenza $W = V i$ la corrente aumenta e la tensione diminuisce, e posso rappresentare la potenza come l'area del rettangolo che si forma tra l'asse delle correnti e l'asse delle tensioni, ad un certo punto la corrente non sarà più lineare quindi avrò non un rettangolo perfetto ma un rettangolo con una curva, questa discrepanza tra la curva e l'eventuale "rettangolo ideale" mi rappresenta la potenza dissipata a causa della resistenza interna.
    \hfill\\\hfill\\\hfill\\\hfill\\
    \subsection{misurazione della f.e.m.}
    Iniziamo col misurare con un Voltmetro la differenza di potenziale ai capi del generatore, e otteniamo un certo valore $V$.\\
    \par
    \begin{wrapfigure}{l}{0.4\textwidth}
        \begin{center}
            \begin{tikzpicture}
                \draw (0,0) to [battery1, invert, l=$\mathcal{E}\quad \pm$,] (0,2);
                \draw (0,2) to [short] (2,2);
                \draw (2,2) to [R, l=$r_x \quad \rho_x$] (2,0);
                \filldraw(2,0) circle (2pt) node [below] {$A$};
                \filldraw(0,0) circle (2pt) node [below] {$B$};
                %voltmetro
                \draw(2,0) to[R, l=$r_x \quad$] (0,0);
                \draw[->] (1.5,0.5) -- (0.5,0.5) node[above, midway] {$R$};
                \draw(2,0) -- (2,-1);
                \draw(0,0) -- (0,-1);
                \draw(0,-1) to [short, voltmeter] (2,-1);
            \end{tikzpicture}
        \end{center}
    \end{wrapfigure}
    \noindent
    Sappiamo che la differenza di potenziale è data sia dalla legge di ohm che dalla relazione con la forza elettromotrice che abbiamo visto prima, quindi avremo che:
    \begin{gather*}
        \Delta V = V_A - V_B = i R = \mathcal{E} - r i
    \end{gather*}
    Se consideriamo quindi una resistenza $R$ in serie alla resistenza interna $r$ del generatore, e risolviamo l'espressione precedente per $i$ troviamo:
    \begin{gather*}
        i = \frac{f_\text{em}}{(R + r)}
    \end{gather*}
    Con la $f_\text{em}$ che ricordiamo la indichiamo proprio con $\E$, in più il fatto che le resistenze si sommino torna poichè sono in serie. Ora ricordando 
    che $V_A - V_B = R I$ e dalla relazione che abbiamo appena scritto possiamo eguagliare questo a:\\
    \begin{gather*}
        R I = \E \frac{R}{R + r} \simeq \E \qquad \text{se $R \gg r$ e $I=0$}
    \end{gather*}
    \subsection{Collegamenti in serie e parallelo}
    Siccome lo abbiamo accennato formalizziamo per bene cosa vuol dire mettere due resistenze in serie e anche cosa vuol dire mettere due resistenze in parallelo.\\
    Per prima cosa ricordiamo che i seguenti circuiti sono equivalenti:
    \begin{center}
        \begin{tikzpicture}[scale = 0.8, transform shape]
            \begin{scope}[yshift= 1.2cm]
                \draw (0,0) to [battery1, invert, l=$\mathcal{E}\quad \pm$,] (0,2);
                \draw (0,2) to [R, l=$R_1$] (2,2);
                \draw (2,2) to [R, l=$R_2$] (2,0);
                \draw (0,0) to[short] (2,0);
            \end{scope}
            \node at(2.7,2) [right] {\huge{$\equiv$}};
            \begin{scope}[xshift=5cm]
                \draw (0,0) to [battery1, invert, l=$\mathcal{E}\quad \pm$,] (0,4);
                \draw (0,4) to [short] (2,4);
                \draw (2,4) to [R, l=$R_1$] (2,2);
                \draw (2,2) to [R, l=$R_2$] (2,0);
                \draw (0,0) to[short] (2,0);   
            \end{scope}        
        \end{tikzpicture}
    \end{center}
    \par
    \begin{wrapfigure}{r}{0.4\textwidth}
        \centering
        \begin{tikzpicture}[scale=0.7, transform shape]
                \draw (0,0) to[R, l=$R_1$] (0,2) to[R, l=$R_2$] (0,4);
                \node at(0.5,2) [right] {\huge{$\equiv$}};
                \draw (2,0) to[R, l=$R_\text{eq}$] (2,4);
                \node at(4.5,2) {Collegamento in serie};
            \end{tikzpicture}
            \hfill\\
        \begin{tikzpicture}[scale=0.7, transform shape]
                \draw(0,0) to[short] (2,0);
                \draw(0,0) to[R, l=$R_1$] (0,2);
                \draw(2,0) to[R, l=$R_2$] (2,2);
                \draw(0,2) to[short] (2,2);
                \draw(1,0) to[short, f=$i$] (1,-0.5);
                \draw(1,2) to[short] (1,2.5);
                \filldraw(1,0) circle (1pt) node[below left] {$n$};
                \filldraw(1,2) circle (1pt) node[above left] {$n$};
                \node at(4.5,1) {Collegamento in parallelo};
            \end{tikzpicture}
    \end{wrapfigure}
    \noindent
    In generale due resistenze sono in serie se sono attraversate dalla stessa corrente, e quindi se sono in serie allora la differenza di potenziale totale è data dalla somma delle differenze di potenziale ai capi di ogni resistenza:
    \begin{gather*}
        R_\text{eq} = R_1 + R_2
    \end{gather*}
    Due resistenze sono in parallelo se sono sottoposte alla stessa differenza di potenziale, e quindi se sono in parallelo allora la corrente totale che attraversa le resistenze è data dalla somma delle correnti che attraversano ogni resistenza:
    \begin{gather*}
        \Delta V = i_1 R_1 = i_2 R_2\\
        i = i_1 + i_2 \\
        i_1 = \frac{\Delta V}{R_1} \qquad i_2 = \frac{\Delta V}{R_2} \\
        i = \frac{\Delta V}{R_1} + \frac{\Delta V}{R_2} = \Delta V \left( \frac{1}{R_1} + \frac{1}{R_2} \right) = \frac{\Delta V}{R_\text{eq}} \\
        \Rightarrow R_\text{eq} = \frac{1}{\frac{1}{R_1} + \frac{1}{R_2}} = \frac{R_1 R_2}{R_1 + R_2}
    \end{gather*}
    \newpage
    \subsection{legge del partitore di tensione}
    \begin{wrapfigure}{r}{0.4\textwidth}
    \centering
        \begin{tikzpicture}
            \draw (0,0) to[R, l=$R_1$] (0,2) to[R, l=$R_2$, f=$i$] (0,4);
            \filldraw (0,0) circle (2pt) node [below] {$A$};
            \filldraw (0,2) circle (2pt) node [left] {$B$};
            \filldraw (0,4) circle (2pt) node [above] {$C$};
        \end{tikzpicture}
    \end{wrapfigure}
    La legge del partitore di tensione dice che se abbiamo due resistenze in serie, allora la differenza di potenziale totale è data dalla somma delle differenze di potenziale ai capi di ogni resistenza, e quindi se vogliamo calcolare la differenza di potenziale ai capi di una resistenza.
    \begin{gather*}
        V_B - V_C = \Delta V_\text{out}\\
        V_A - V_C = \Delta V \qquad V_A - V_C = iR_1 + iR_2 \\
        \Rightarrow V_A - V_C = i (R_1 + R_2) \\
    \end{gather*}
    e quindi anche:
    \begin{gather*}
        V_A - V_C = \Delta V_\text{in} \\
        V_B - V_C = i R_2\\
    \end{gather*}
    Esplicitiamo $i$ dalla prima equazione e sostituiamolo nella seconda:
    \begin{gather*}
        V_B - V_C = \frac{V_A - V_C}{R_1 + R_2} R_2 = \Delta V_\text{in} \frac{R_2}{R_1 + R_2}
    \end{gather*}
    Siccome $V_B - V_C = \Delta V_0$
    \begin{gather*}
        \Delta V_\text{out} = \Delta V_\text{in} \frac{R_2}{R_1 + R_2}
    \end{gather*}
    \subsection{Misurazione fem}
    \begin{wrapfigure}{l}{0.4\textwidth}
            \begin{tikzpicture}
                \draw (0,0) to [battery1, invert, l=$\E_x\quad \pm$,] (0,2);
                \draw (0,2) to [short, f=$i$] (2,2);
                \draw (2,2) to [R, l=$\rho_x$] (2,0);
                \filldraw(2,0) circle (2pt) node [below, right] {$A$};
                \filldraw(0,0) circle (2pt) node [below, left] {$B$};
                \draw(2,0) to [short, R, l=$R$, f=$i$] (0,0);
                \draw(0,0) to [short] (0,-2);
                \draw(0,-2) to [short, voltmeter] (2,-2);
                \draw(2,-2) to [short] (2,0);
                \draw(2,-2) to (2,-4);
                \draw(0,-2) to (0,-4);
                \draw(2,-4) to [short, R, l=$R$] (0,-4);
                \draw(2,-4) to [short] (2,-6);
                \draw(0,-4) to [short] (0,-6);
                \draw(0,-6) to [short, voltmeter] (2,-6);
        \end{tikzpicture}
    \end{wrapfigure}
    Questo sarà il circuito che utilizzeremo anche nella prima esperienza e serve per la prima parte ad avere una buona stima della f.e.m. che vogliamo misurare.\\
    Il nostro scopo è ricavare le incognite $\E_x$ ed $\rho_x$, che ricordiamo essere rispettivamente la forza elettromotrice del generatore e la sua resistenza interna.\\
    Per ricavarle effettuiamo due misurazioni, sul circuito illustrato, e posizioniamo delle resistenze in modo da ottenere il seguente sistema.
    \begin{gather*}
        \begin{cases}
            (R + 2\rho_x) V_2 = R \E_x\\
            (R + \rho_x) V_1 = R \E_x
        \end{cases}
    \end{gather*}
    Dove $V_1$ e $V_2$ sono le due misurazioni che otteniamo dai voltmetri, e $R$ è la resistenza che abbiamo posizionato.\\
    Risolvendo il sistema troviamo:
    \begin{gather*}
        \E_x = \frac{(R + 2\rho_x) V_2}{R} = \frac{(R + \rho_x) V_1}{R} \\
        \Rightarrow (R + 2\rho_x) V_2 = (R + \rho_x) V_1 \\
        \Rightarrow R (V_2 - V_1) = \rho_x (V_1 - 2 V_2) \\
        \Rightarrow \rho_x = R \frac{V_2 - V_1}{V_1 - 2 V_2}
    \end{gather*}
    E sostituiendo $\rho_x$ in una delle due equazioni del sistema troviamo:
    \begin{gather*}
        \E_x = \frac{V_1 V_2}{2V_2 -V_1}
    \end{gather*}
    \hfill\\\hfill\\
    \subsection{un tipo di multimetro}
    Caratteristiche generali del multimetro fluke 114:
        \begin{enumerate}[1)]
            \item in DC V / DC mV $\qquad 0,5\% + 2$ digit. (sull'ultima misurata)
            \item $10 k \Omega - 40 M \Omega \qquad 1,5\% + 2$ digit. (sull'ultima misurata)
        \end{enumerate}
        \begin{example}
            \begin{gather*}
                R= 11,09 M\Omega\\
                \Delta R = 1,5 \% \cdot 11,09 M \Omega + 0,02 M \Omega \sim 0,17 M \Omega + 0,02 M \Omega\\
                R = (11,09 \pm 0,17) M \Omega
            \end{gather*}
        \end{example}
\newpage\noindent
\section{4/3/26}
\subsection{Prima esperienza}
\leavevmode
\begin{wrapfigure}{r}{0.4\textwidth}
    \centering
    \begin{tikzpicture}
        \draw (0,0) to [generic, invert, l=$\E \ \pm$] (0,1);
        \draw (0,0) to[short] (0,-0.5);
        \draw (0,1) to[short] (0,1.5);
        \draw (0,-0.5) to[short] (2,-0.5);
        \draw (0,1.5) to[short] (2,1.5);
        \draw (2,1.5) to[short, f=$i$] (2,-0.5);
        \draw[<-] (-1,0) -- (-1,1) node[left, at start] {$\vv{E}_e$};
        \node at(1,1) [above] {$\Gamma$};
        \filldraw(0,-0.1) circle (1pt) node[left] {$B$};
        \filldraw(0,1.1) circle (1pt) node[left] {$A$};
        \draw[->] (0.1,0.2) -- (0.1,0.8) node[left, midway] {$i$};
        \draw[->] (0.5,0.1) -- (0.5,1) node[right, midway] {$\vv{E}_s'$};
    \end{tikzpicture}
    \caption{Schema del generatore}
    \end{wrapfigure}\noindent
Sappiamo che all'equilibrio si ha:
\begin{gather*}
    \vv{E}_e = -\vv{E}'_s
\end{gather*}
Dove si ricorda che $\vv{E}'_s$ è il campo elettrostatico all'interno del generatore, e che è in contrapposizione con il campo $\vv{E}_e$ che rappresenta la forza che separa le cariche all'interno del generatore. All'interno del 
generatore entrambi i campi sono conservativi.\\
Da questi possiamo calcolare la f.e.m. che è data da:
\begin{gather*}
    f.e.m. = \E = \int_{A}^{B} \vv{E}_e \cdot \vv{dl} = \int_{A}^{B} -\vv{E}'_s \cdot \vv{dl}\\
    = \int_{A}^{B} \vv{E}'_s \cdot \vv{dl} = V_A-V_B
\end{gather*}
Di conseguenza possiamo avere la definizione operativa di f.e.m. come la differenza di potenziale che si misura ai capi del generatore quando non c'è corrente che lo attraversa. Ed è ciò che dobbiamo riuscire a misurare\\
\begin{gather*}
    fem = f = \mathcal{E} = (V_A - V_B)_\text{misurata} \vert_{i=0}
\end{gather*}
E si misura in volt $[V]$
In generale consideriamo una situazione rappresetata dal circuito sulla sinistra, sulla destra è anche riportato il grafico che rappresenta la decaduta lineare della resistenza fino ad un certo valore:
\begin{center}
    \begin{tikzpicture}
        \draw (0,0) to [battery1, invert, l=$f\quad \pm$,] (0,2);
        \draw (0,2) to [short] (2,2);
        \draw (2,2) to [R, l=$r$] (2,0);
        \filldraw(2,0) circle (2pt) node [below] {$A$};
        \filldraw(0,0) circle (2pt) node [below] {$B$};
        \draw(2,0) to [short] (0,0);
        \draw[->] (1.5,1.8) -- (1.5,0.2) node[right, midway] {$i$};
    \end{tikzpicture}
    \begin{tikzpicture}
            %assi x y
            \draw[->] (0,0) -- (3,0) node[below] {$i$};
            \draw[->] (0,0) -- (0,3) node[left] {$V_A-V_B$};
            %linea
            \draw[thick] (0,2) -- (2,0.5) node[midway, above] {$r$} node[at start, left]{$f$};
            \draw[thick, cyan] (2,0.5) to[bend left = 10] (2.5,0);
            \node at(2.5,0) [below] {$i_m$};
    \end{tikzpicture}
\end{center}
Come già spiegato precedentemente il potenziale tra $A$ e $B$ è dato da:
\begin{gather*}
    (V_A - V_B) = f - r i\\
    = \mathcal{E} -i \rho
\end{gather*}
Questo ovviamente vale fino a che ci manteniamo nel regime di proporzionalità, ovvero non siamo vicini alla corrente massima $i_m$ che il generatore riesce a sostenere.\\
    \newpage\noindent
    Riprendiamo ora i circuiti della scorsa lezione:
    \begin{center}
        \begin{tikzpicture}
            \draw (0,0) to [battery1, invert, l=$f\quad \pm$,] (0,2);
            \draw (0,2) to [short] (2,2);
            \draw (2,2) to [R, l=$r$] (2,0);
            \filldraw(2,0) circle (2pt) node [below, right] {$A$};
            \filldraw(0,0) circle (2pt) node [below, left] {$B$};
            \draw(2,0) to [short, R, l=$R$] (0,0);
            \draw[->] (1.5,1.8) -- (1.5,0.2) node[right, midway] {$i$};
            \draw(0,0) to [short] (0,-2);
            \draw(0,-2) to [short, voltmeter] (2,-2);
            \draw(2,-2) to [short] (2,0);
        \end{tikzpicture}
        \begin{tikzpicture}
            \draw (0,0) to [battery1, invert, l=$f\quad \pm$,] (0,2);
            \draw (0,2) to [short] (2,2);
            \draw (2,2) to [R, l=$r$] (2,0);
            \filldraw(2,0) circle (2pt) node [below, right] {$A$};
            \filldraw(0,0) circle (2pt) node [below, left] {$B$};
            \draw(2,0) to [short, R, l=$R$] (0,0);
            \draw[->] (1.5,1.8) -- (1.5,0.2) node[right, midway] {$i$};
            \draw(0,0) to [short] (0,-2);
            \draw(0,-2) to [short, voltmeter] (2,-2);
            \draw(2,-2) to [short] (2,0);
            \draw(2,-2) to (2,-4);
            \draw(0,-2) to (0,-4);
            \draw(2,-4) to [short, R, l=$R$] (0,-4);
            \draw(2,-4) to [short] (2,-6);
            \draw(0,-4) to [short] (0,-6);
            \draw(0,-6) to [short, voltmeter] (2,-6);
        \end{tikzpicture}
    \end{center}
    Eravamo arrivati a ricavare la differenza di potenziale in base alla misurazione dei multimetri si aveva dal primo circuito:
    \begin{gather*}
        (V_A - V_B)_1 = V_1
    \end{gather*}
    e dal secondo circuito:
    \begin{gather*}
        (V_A - V_B)_2 = V_2
    \end{gather*}
    Questo metodo di fatto ci permette di ricavare la fem anche se non siamo a corrente nulla, e per quanto queste leggi siano giuste e valide, per definizione questa non è esattamente la fem, ma di sicuro una suoa buona stima, questo metodo semplicemente non può essere scalato 
    con strumenti più precisi per ottenere risultati molto migliori, di sicuro però mi è utile sia per stimare la fem sia per misurare abbastanza accuratamente la resistenza interna del generatore.
\subsection{Metodo del divisore di tensione}
Per questo metodo dobbiamo introdurre un nuovo oggetto, prendiamo in considerazione il seguente circuito:
\begin{center}
    \begin{tikzpicture}
        \draw (0,-2) to[battery1, invert, l=$f \quad \pm$,] (0,2);
        \draw (0,2) to [short] (2,2);
        \draw (2,2) to[vR, l=$R_v$] (2,0) to [vR, l=$R_v$] (2,-2);
        \draw (2,-2) to [short] (0,-2);
        \filldraw(2,-2) circle (2pt);
        \filldraw(2,2) circle (2pt);
        \draw(1.5,-2.3) rectangle (2.5,2.3);
        \draw(2,0) to[ammeter, f=$i{=}0$] (4,0) to[R, l=$\rho_x$] (6,0);
        \draw(6,0) to[battery1, invert, l=$\pm\quad \E_x$,] (6,-2);
        \draw(6,-2) to [short] (0,-2);
    \end{tikzpicture}
\end{center}
Il rettangolo rappresentato nel circuito sovrastante è proprio un divisore di tensione che è un dispositivo che ha resistenze variabili che ci permette di cambiare la tensione in uscita in base ad un cursore che muoviamo noi.\\
Di fatto come facciamo a misurare la fem del generatore $\E_x$?\\
Intanto si deve sapere che non ci interessa la fem del generatore $f$ che nel circuito è rapresentato a destra, questo perchè giocando col cursore del divisore di tensione 
possiamo ottenere il giusto rapporto di resistenze fino a che l'amperometro (che si ricorda essere uno strumento che misura proprio la corrente) ci dirà che nella maglia di destra non passerà più corrente, e quindi la tensione che misureremo ai capi del generatore $\E_x$ sarà proprio la sua fem.\\
Questo è chiamato \textbf{metodo potenziometrico} e in laboratorio sarà composto da due step o circuiti:\\
\begin{center}
\begin{tikzpicture}[scale=0.7, transform shape]
        \draw (0,-2) to[battery1, invert, l=$f \quad \pm$,] (0,2);
        \draw (0,2) to [short] (2,2);
        \draw (2,2) to[vR, l=$R_v$] (2,0) to [vR, l=$R_v$] (2,-2);
        \draw (2,-2) to [short] (0,-2);
        \filldraw(2,-2) circle (2pt);
        \filldraw(2,2) circle (2pt);
        \draw(1.5,-2.3) rectangle (2.5,2.3);
        \draw(2,0) to[ammeter, f=$i{=}0$] (4,0) to[R, l=$r_x$] (6,0);
        \draw(6,0) to[battery1, invert, l=$\pm\quad \E_x$,] (6,-2);
        \draw(6,-2) to [short] (0,-2);
    \end{tikzpicture}
    \begin{tikzpicture}[scale=0.7, transform shape]
        \draw (0,-2) to[battery1, invert, l=$f \quad \pm$,] (0,2);
        \draw (0,2) to [short] (2,2);
        \draw (2,2) to[vR, l=$R_v$] (2,0) to [vR, l=$R_v$] (2,-2);
        \draw (2,-2) to [short] (0,-2);
        \filldraw(2,-2) circle (2pt);
        \filldraw(2,2) circle (2pt);
        \draw(1.5,-2.3) rectangle (2.5,2.3);
        \draw(2,0) to[ammeter, f=$i{=}0$] (4,0) to[R, l=$r_R$] (6,0);
        \draw(6,0) to[battery1, invert, l=$\pm\quad \E_R$,] (6,-2);
        \draw(6,-2) to [short] (0,-2);
    \end{tikzpicture}
\end{center}
Si calcolano le due f.e.m.  di entrambi i generatori $\E_x$ e $\E_R$ con il metodo potenziometrico, si eseguono due misure proprio per ricavare le resistenze interne $r_x$ e $r_R$ dei generatori.\\
Con questo metodo si riescono a raggiungere discrete precisioni fino a: $\frac{\Delta V}{V} \sim 10^{-9}$per le tensioni e $\frac{\Delta R}{R} \sim 10^{-9}$ per le resistenze.
\subsection{leggi nodi}
Questa legge ci dice che la somma algebrica delle correnti che entrano in un nodo è uguale alla somma algebrica delle correnti che escono da quel nodo, e quindi se consideriamo un nodo con $n$ correnti che entrano e $m$ correnti che escono, allora avremo:
\begin{gather*}
    \sum_k i_k = 0 \qquad \nabla \cdot \vv{J} = 0\\
    \int_S \vv{J} \cdot \vv{dS} = 0
\end{gather*}
\subsection{Legge delle maglie}
Questa legge ci dice che la somma algebrica delle differenze di potenziale lungo una maglia chiusa è uguale a zero, e quindi se consideriamo una maglia con $n$ differenze di potenziale, allora avremo:
\begin{gather*}
    \sum_k \Delta V_k = 0 \qquad \nabla \land \vv{E} = 0\\
    \oint_\Gamma \vv{E} \cdot \vv{dl} = 0
\end{gather*}
Ricordiamo che noi schematizziamo i circuiti come se le resistenze fossero puntiformi, ma la legge che vale punto per punto di fatto è:
\begin{gather*}
    \vv{E}(\vv{r}) = \rho \vv{J}(\vv{r})
\end{gather*}
\subsection{esercizi con legge di ohm generalizzata}
Prendiamo il circuito:
\begin{center}
    \begin{tikzpicture}
        \draw(0,0) to[R, l=$R$] (2,0) to[battery1, invert, l=$- \ +$] (3,0) to[R, l=$R$] (5,0) to[battery1, invert, l=$- \ +$] (6,0) to[R, l=$R$] (8,0);
        \filldraw(0,0) circle (2pt) node [below] {$A$};
        \filldraw(8,0) circle (2pt) node [below] {$B$};
    \end{tikzpicture}
\end{center}
Abbiamo che:
\begin{gather*}
    \int_{A}^{B} \vv{E} \cdot \vv{dl} = \int_{A}^{B} \rho_R \vv{J} \cdot \vv{dl}
\end{gather*}
Ricordando che $\vv{E} = \vv{E}_e + \vv{E}_s$ in cui $\vv{E}_e$ è la parte del campo elettrico conservativo e $\vv{E}_s$ è la parte del campo elettrico non conservativo, allora possiamo scrivere:
\begin{gather*}
    \int_{A}^{B} \vv{E}_e \cdot \vv{dl} + \int_{A}^{B} \vv{E}_s \cdot \vv{dl} = \int_{A}^{B} \rho_R \vv{J} \cdot \vv{dl}
\end{gather*}
Ora considero che la parte del campo elettrico conservativo è proprio la differenza di potenziale tra il punto $A$ e il punto $B$, e che il campo non conservativo lo posso dividere nei seguenti:
\begin{gather*}
    (V_A - V_B) +  \int_{-}^{+} \vv{E}_1 \cdot \vv{dl} + \int_{+}^{-} \vv{E}_2 \cdot \vv{dl} = \int_{A}^{B} \rho_R \vv{J} \cdot \vv{dl}
\end{gather*}
Ora ho che $\int_{-}^{+} \vv{E}_1 \cdot \vv{dl} = f_1$ e $\int_{+}^{-} \vv{E}_2 \cdot \vv{dl} = -f_2$, in più le resistenze che ho: $R_1$, $R_2$, $R_3$ le posso sommare in base alla loro corrente, quindi posso scrivere:
\begin{gather*}
    (V_A - V_B) + f_1 - f_2 = R_1 i_1 + R_2 i_2 + R_3 i_3 = (R_1 + R_2 + R_3)i
\end{gather*}
\subsection{Seconda legge di kirckoff applicata alla maglia}
\begin{gather*}
    \sum_\text{maglia} f = (\sum_\text{maglia} R i)
\end{gather*}
Questo vuol dire che la somma delle f.e.m. che attraversano una maglia è uguale alla somma delle cadute di tensione che attraversano quella maglia, e quindi se consideriamo una maglia con $n$ f.e.m. e $m$ resistenze, allora avremo:
\begin{gather*}
    \sum_{k=1}^{n} f_k = \sum_{k=1}^{m} R_k i_k
\end{gather*}
\newpage
\section{9/3/26}
La scorsa volta abbiamo impostato il seguente esercizio:
\begin{center}
    \begin{tikzpicture}
        \draw (0,0) to [battery1, invert, l=$\mathcal{E}_1\quad \pm$,] (0,3);
        \draw (0,3) to [R, l=$R_1$] (2,3);
        \draw (2,0) to [battery1, l=$\mathcal{E}_2 \quad \mp$] (2,1);
        \draw (2,1) to[ R, l=$R_2$] (2,3);
        \draw(2,0) to [short] (0,0);
        \draw(2,0) to[short] (4,0);
        \draw(2,3) to[short] (4,3);
        \draw(4,0) to[R, l=$R_3$] (4,3);
        \node[scale = 2.5] at(1,1.5) {$\circlearrowright$};
        \node[scale = 2.5] at(3,1.5) {$\circlearrowright$};
        \node at(1,1.5) {$I_1$};
        \node at(3,1.5) {$I_2$};
    \end{tikzpicture}
\end{center}
Procediamo a risolverlo con il metoodo delle maglie che abbiamo introdotto la scorsa volta, consideriamo le correnti di maglia $I_1$ e $I_2$.\\
Abbiamo che per $I_1$ stiamo percorrendo i due generatori da $-$ a $+$ quindi la corrente di maglia sarà $\E_1 + \E_2$, Per $I_2$ invece stiamo percorrendo $\E_2$ da $-$ a $+$ quindi la corrende di maglia sarà $-\E_2$.\\
Per le resistenze ci basta sommare in base alle correnti di maglia per la 1: $I_1 R_1$ (poichè solo della magl. 1) più $(I_1-I_2) R_2$ poichè in comune e consideriamo $I_1$ positiva poichè è quella della maglia che stiamo considerando, quindi per la maglia 2: $I_2 R_3$ più $(I_2 -I_1) R_2$.
\begin{gather*}
    \begin{cases}
        \mathcal{E}_1 + \mathcal{E}_2 = I_1 R_1 + (I_1 - I_2) R_2\\
        -\mathcal{E}_2 = I_2 R_3 + (I_2 - I_1) R_2
    \end{cases}
\end{gather*}
Isoliamo i termini in $I_1$ e $I_2$:
\begin{gather*}
    \begin{cases}
        \mathcal{E}_1 + \mathcal{E}_2 = I_1 (R_1 + R_2) - I_2 R_2\\
        -\mathcal{E}_2 = I_2 (R_2 + R_3) - I_1 R_2
    \end{cases}
\end{gather*}
Da quì possiamo sviluppare il prblema in forma matriciale:
\begin{gather*}
    \begin{pmatrix}
        \mathcal{E}_1 + \mathcal{E}_2\\
        -\mathcal{E}_2
    \end{pmatrix} = \underbrace{\begin{pmatrix}
        R_1 + R_2 & -R_2\\
        -R_2 & R_2 + R_3
    \end{pmatrix}}_{R_{kl}} \cdot
    \underbrace{\begin{pmatrix}
        I_1\\
        I_2
    \end{pmatrix}}_{I_e}
\end{gather*}
Da qui otteniamo che $\mathcal{E}_k = R_{kl} I_e$. Quindi una matrice $2\times2$ per una matrice $2\times1$:\\
\subsection{Metodo di Cramer}
Introduciamo velocemente il metodo di Cramer per risolvere sistemi lineari:
Preso un sistema e la sua forma amtriciale:
\begin{gather*}
    \begin{cases}
        R_{11} I_1 + R_{12} I_2 = \mathcal{E}_1\\
        R_{21} I_1 + R_{22} I_2 = \mathcal{E}_2
    \end{cases}
    \Rightarrow
    \begin{pmatrix}
        \mathcal{E}_1\\
        \mathcal{E}_2
    \end{pmatrix} = \begin{pmatrix}
        R_{11} & R_{12}\\
        R_{21} & R_{22}
    \end{pmatrix} \cdot
    \begin{pmatrix}
        I_1\\
        I_2
    \end{pmatrix}
\end{gather*}
Si calcolano i determinanti delle seguenti matrici:
\begin{gather*}
    det(R) = \left\lvert \begin{matrix}
        R_{11} & R_{12}\\
        R_{21} & R_{22}
    \end{matrix} \right\rvert \qquad
    det(R_1) = \left\lvert \begin{matrix}
        \mathcal{E}_1 & R_{12}\\
        \mathcal{E}_2 & R_{22}
    \end{matrix} \right\rvert \qquad
    det(R_2) = \left\lvert \begin{matrix}
        R_{11} & \mathcal{E}_1\\
        R_{21} & \mathcal{E}_2
    \end{matrix} \right\rvert
\end{gather*}
Dove la colonna sostituita corrisponde alla corrente $I_k$-esima che vogliamo calcolare, nel caso delle 2x2 la $1^a$ colonna per $I_1$ e la $2^a$ colonna per $I_2$.\\
Se $det(R) \neq 0$ allora il sistema è risolvibile (cosa che nel nostro caso è sempre vera poiche le resistenze sono $>0$) e le soluzioni sono:
\begin{gather*}
    I_1 = \frac{det(R_1)}{det(R)} \qquad I_2 = \frac{det(R_2)}{det(R)}
\end{gather*}
Questo è un sistema lineare ed è risolvibile con il metodo di Cramer, e quindi otteniamo:
\begin{gather*}
    i_3 = I_2 = \frac{\left\lvert \begin{matrix}
        R_1 + R_2 & \mathcal{E}_1 + \mathcal{E}_2\\
        -R_2 & -\mathcal{E}_2
    \end{matrix} \right\rvert }{\left\lvert \begin{matrix}
        R_1 + R_2 & -R_2\\
        -R_2 & R_2 + R_3
    \end{matrix} \right\rvert }
    \Rightarrow I_1 = \frac{\left\lvert \begin{matrix}
        \mathcal{E}_1 + \mathcal{E}_2 & -R_2\\
        -\mathcal{E}_2 & R_2 + R_3
    \end{matrix} \right\rvert }{\left\lvert \begin{matrix}
        R_1 + R_2 & -R_2\\
        -R_2 & R_2 + R_3
    \end{matrix} \right\rvert }
\end{gather*}
Si ricorda che nel metodo delle maglie le correnti di maglia sono indicate con $I$ e quelle di ramo con $i$.\\ 
Svolgendo i conti dei determinanti arriviamo all'espressione:
\begin{gather*}
    I_2 = \frac{-\mathcal{E}_2(R_1 + R_2) + (\mathcal{E}_1 + \mathcal{E}_2)R_2}{(R_1 + R_2)(R_2 + R_3) - R_2^2}\\
    = \frac{-\mathcal{E}_2 R_1 - \cancel{\mathcal{E}_2 R_2} + \mathcal{E}_1 R_2 + \cancel{\mathcal{E}_2 R_2}}{R_1 R_2 + R_1 R_3 + \cancel{R_2^2} + R_2 R_3 - \cancel{R_2^2}}\\
    I_2 = i_3 = \frac{-\mathcal{E}_2 R_1 + \mathcal{E}_1 R_2}{S}
\end{gather*}
Consideriamo ora di lavorare non con due maglie ma con $n$ maglie, e quindi con $n$ correnti di maglia, e quindi con $n$ equazioni, in questo caso avremo che:
\begin{gather*}
    R_{kl} = \begin{pmatrix}
        R_{11} & R_{12} & \dots & R_{1n}\\
        R_{21} & R_{22} & \dots & R_{2n}\\
        \vdots & \vdots & \ddots & \vdots\\
        R_{n1} & R_{n2} & \dots & R_{nn}
    \end{pmatrix}
\end{gather*}
Cerco ora la corrente $I_k$:
\begin{gather*}
    I_k = \frac{\left\lvert \begin{matrix}
        R_{11} & \dots & R_{1k-1} & \E_1 & R_{1k+1} & \dots & R_{1n}\\
        R_{21} & \dots & R_{2k-1} & \E_2 & R_{2k+1} & \dots & R_{2n}\\
        \vdots &  & \vdots & \vdots & \vdots &  & \vdots\\
        R_{n1} & \dots & R_{nk-1} & \mathcal{E}_n & R_{nk+1} & \dots & R_{nn}
    \end{matrix} \right\rvert }{\left\lVert R_{kl}\right\rVert }
\end{gather*}
Stiamo usando la notazione di $\left\lVert \cdot \right\rVert$ che rappresenta il determinante della matrice.\\
Se applico gli sviluppi di laplace alla matrice al numeratore, posso riscrivere questa espressione come una somma di contributi:
\begin{gather*}
    I_k = \mathcal{E}_1 \frac{\left\lVert C_{1k} \right\rVert }{\left\lVert R_{kl}\right\rVert } + \dots + \mathcal{E}_n \frac{\left\lVert C_{nk} \right\rVert }{\left\lVert R_{kl} \right\rVert }
\end{gather*}
Dove $C_{ik} = (-1)^{i+k} \ \det(M_{ik})$, dove $M_{ik}$ sono le matrici che si ottengono eliminando la $i$-esima riga e la $k$-esima colonna dalla matrice $R_{kl}$.\\ 
Si trova che $I_k = \mathcal{E}_k \frac{\left\lVert C_{kk} \right\rVert }{\left\lVert R_{kl} \right\rVert } = \frac{\mathcal{E}_k}{R_T}$ quindi si è trovato che dipende da una resistenza che possiamo esplicitare in questo modo:
\begin{gather*}
    R_T = \frac{\left\lVert R_{kl} \right\rVert }{\left\lVert C_{kk} \right\rVert }
\end{gather*}
Si noti che anche sisccome i segni di $C_{kk}$ si alternano potrebbe sembrare che possa venire una resistenza negativa, ma siccome siamo su indici della diagonale questo non accadrà mai poichè si ricorda che quegli indici daranno sempre coefficienti positivi: $(-1)^{k+k} = (-1)^{2k} = 1$
\tiny{$C$ a lezione non è stato presentato, si parlò solo di $M$ e il prof disse i minori, i segni non tornerebbero ma bho} \normalsize
\subsection{Teorema di Thevenin}
È un teorema che afferma che un qualsiasi rete elettrica lineare può essere espresso come un solo generatore ideale di tensione $\E_{th}$ 
in serie ad un unico resistore $R_{th}$. Con un circuito lineare si intende un circuito che è composto da generatori e resistenze che hanno 
caratteristiche costanti, in generale anche altri componenti ma che non variano al variare del segnale elettrico (tensione e corrente).
\begin{center}
    \begin{tikzpicture}
    \begin{scope}
            \draw(-0.5,0) rectangle (2,2);
            \node at(0.7,1.4){\small{RETE}};
            \node at(0.5,0.6){\small{LINEARE}};
            \draw (1.5,0.5) -- (2.5,0.5) node[below] {$A$};
            \draw (1.5,1.5) -- (2.5,1.5) node[above] {$B$};\
            \filldraw (2.5,0.5) circle (2pt);
            \filldraw (2.5,1.5) circle (2pt);
    \end{scope}
    \node at(4,1) {$\Huge{\equiv}$};
    \begin{scope}[shift={(7,0)}]
            \draw (0,0) to [battery1, invert, l=$\mathcal{E}_\text{th}\quad \pm$,] (0,2);
            \draw (0,2) to[R, l=$R_\text{th}$] (2,2);
            \draw (0,0) to[short] (2,0);
            \filldraw(2,0) circle (2pt) node [below] {$A$};
            \filldraw(2,2) circle (2pt) node [above] {$B$};
    \end{scope}
    \end{tikzpicture}
\end{center}
\newpage
\begin{example}
    Vediamo un caso applicativo di questo teorema, vediamo l'equivalenza della stella triangolo, di fatto dimostreremo che i seguenti circuiti sono equivalenti:
    \begin{center}
        \begin{tikzpicture}
            \begin{scope}[xshift=-3cm,yshift=-1cm]
                %\draw (0,0) to [battery1, invert, l=$\mathcal{E}_1\quad \pm$,] (0,4);
                %\draw(0,4) to [short] (3,4);
                %\draw(0,0) to[short] (3,0);
                %\filldraw (3,4) circle (2pt) node [above] {$A$};
                %\filldraw (3,0) circle (2pt) node [above] {$B$};
                \draw (1.5,2) to[R, l=$R_{AB}$, invert] (3,4);
                \draw (3,4) to[R, l=$R_{AC}$] (4.5,2);
                \draw (1.5,2) to[R, l=$R_{BC}$] (4.5,2);
                \filldraw (3,4) circle (2pt) node [above] {$A$};
                \filldraw (1.5,2) circle (2pt) node [above] {$B$};
                \filldraw (4.5,2) circle (2pt) node [above] {$C$};
                %\draw (4.5,2) to[R, l=$R_3$] (3,0);
                %\draw (1.5,2) to[R, l=$R_4$] (3,0);
            \end{scope}
            \node[scale=1.5] at(2.5,2) {$\equiv$};
            \begin{scope}[xshift=2cm]
                \draw (1.5,0) to[R, l=$R_B$, invert] (3,2);
                \draw (3,2) to[R, l=$R_C$] (4.5,0);
                \draw (3,2) to[R, l=$R_A$] (3,4);
                \filldraw (3,4) circle (2pt) node[above] {$A$};
                \filldraw (1.5,0) circle (2pt) node[above] {$B$};
                \filldraw (4.5,0) circle (2pt) node[above] {$C$};
            \end{scope}
        \end{tikzpicture}
    \end{center}
Guardiamo segmento per segmento della stella a cosa corrispondono le resistenze del triangolo:
\begin{gather*}
    \text{segmento AB} \qquad R_A + R_B = R_{AB} \parallel (R_{AC} + R_{BC}) = \frac{R_{AB}(R_{AC} + R_{BC})}{R_{AB} + R_{AC} + R_{BC}}\\
    \text{segmento AC} \qquad R_A + R_C = R_{AC} \parallel (R_{AB} + R_{BC}) = \frac{R_{AC}(R_{AB} + R_{BC})}{R_{AB} + R_{AC} + R_{BC}}\\
    \text{segmento BC} \qquad R_B + R_C = R_{BC} \parallel (R_{AB} + R_{AC}) = \frac{R_{BC}(R_{AB} + R_{AC})}{R_{AB} + R_{AC} + R_{BC}}
\end{gather*}
Sommando tutto otteniamo:
\begin{gather*}
    2R_A + 2R_B + 2R_C = \frac{R_{AC}(R_{AB} + R_{BC})}{\Sigma R} + \frac{R_{BC}(R_{AB} + R_{AC})}{\Sigma R} + \frac{R_{AB}(R_{AC} + R_{BC})}{\Sigma R}\\
    2(R_A + R_B + R_C) = \frac{2(R_{AB}R_{AC} + R_{AB}R_{BC} + R_{AC}R_{BC})}{\Sigma R}
\end{gather*}
Dividendo per due otteniamo la somma delle resistenze della stella:
\begin{gather*}
    R_A + R_B + R_C = \frac{R_{AB}R_{AC} + R_{AB}R_{BC} + R_{AC}R_{BC}}{\Sigma R}
\end{gather*}
Da questa che è la somma delle tre resistenze togliamo $R_B + R_C$ che è espressa proprio nella terza equazione quando abbiamo impostato l'equivalenza nel ramo $BC$ all'inizio.\\
Svolgendo i conti otteniamo che:
\begin{gather*}
    R_A = \frac{R_{AB}R_{AC} + \cancel{R_{AB}R_{BC}} + \cancel{R_{AC}R_{BC}} - (\cancel{R_{BC}R_{AB}} + \cancel{R_{BC}R_{AC}})}{\Sigma R}
\end{gather*}
Il valore della resistenza di stella è dato dal prodotto delle resistenze del triangolo concorrenti nel nodo, diviso la somma totale:
\begin{gather*}
    R_A = \frac{R_{AB} R_{AC}}{R_{AB} + R_{AC} + R_{BC}}
\end{gather*}
Analogamente si ricavano $R_B$ e $R_C$ per permutazione degli indici.\\
Essendo arrivati proprio all'equazione di $R_A$ possiamo dire che presi i valori delle resistenze del triangolo, qualunque essi siano, possiamo sempre isolare le resistenze per descrivere in maniera diversa i vari collegamenti tra $A$ $B$ e $C$.\\
Di fatto il risultato che abbiamo ottenuto è che preso un sistema di quel tipo abbiamo visto che è sempre risolubile, anche se sembra un risultato banale ricordiamoci che per thevenin qualunque circuito possiamo scrivere come sistema lineare avrà sempre soluzione.
\end{example}
\newpage
\leavevmode
\begin{wrapfigure}{r}{0.4\textwidth}   
        \begin{tikzpicture}
            \draw (0,0) to [battery1, invert, l=$\mathcal{E}_\text{th}\quad \pm$,] (0,2);
            \draw (0,2) to[R, l=$R_\text{th}$] (2,2);
            \draw (0,0) to[short] (2,0);
            \filldraw(2,0) circle (2pt) node [below] {$A$};
            \filldraw(2,2) circle (2pt) node [above] {$B$};
        \end{tikzpicture}
    \caption{Circuito equivalente}
\end{wrapfigure}
\begin{proof}[Dimostrazione del th. di Thevenin]
        Ora che abbiamo fatto un esempio incentriamoci ora sulla dimostrazione del teorema di Thevenin.\\
        Iniziamo col ribadire il concetto, thevenin enuncia che preso un circuito lineare e due morsetti $A$ e $B$ di questo, se ne misuriamo la differenza di potenziale, poi cortocircuitiamo tutti i generatori e calcoliamo la resistenza complessiva, possiamo riscrivere il circuito con un singolo generatore di f.e.m. $\E_{th}$ uguale alla tensione della differenza di potenziale misurata e una unica resistenza $R_{th}$ pari al calcolo complessivo di tutte le resistenze.\\
        Cortocircuitare i geneartori vuol dire semplicemente "rimuoverli" e considerarli come semplici conduttori (di fatto nel circuito si sostituiscono con un filo).\\
        Come ipotesi abbiamo quindi che i generatori devono cortocircuitare.\\
        Ricordiamo che la fem a corrente nulla è proprio la differenza di potenziale
        \begin{gather*}
            \E_{th} = V_A - V_B |_{i=0}\\
        \end{gather*}
        In questo caso siamo a circuito aperto quindi a corrente nulla e a tensione massima.
        Dalla linearità di decadimento della resistenza ho:
        \begin{gather*}
            V_{AB} = \E_{th} - R_{th} * I_{cc}
        \end{gather*}
        Dove $I_{cc}$ è la corrente di cortocircuito, ovvero la corrente che scorrerebbe tra $A$ e $B$ se li cortocircuitassimo, e quindi se $I_{cc} = 0$ allora $V_{AB} = \E_{th}$, ma se invece $I_{cc} \neq 0$ allora $V_{AB} \neq \E_{th}$,
        Ma siccome sono a circuito chiuso(i morsetti $A$ e $B$ sono in corto circuito) e ho che la corrente ha valore massimo e vale proprio $I_{cc}$ ho che $V_{AB} = 0$\\
        Quindi posso riscrivere la formula come:
        \begin{gather*}
            0 = \E_{th} - R_{th} * I_{cc}\\
            R_{th} = \frac{\E_{th}}{I_{cc}}
        \end{gather*}
    \end{proof}
\newpage
\section{11/3/26}
\subsection{galvanometro}
È lo strumento base con cui si fanno le varie misure in laboratorio.\\
Di base è un misuratore di corrente, o meglio è un misuratore di corrente che è in grado di misurare l'intensità di correnti molto piccole, dell'ordine dei nanoampere, e quindi è uno strumento molto sensibile.\\ 
\begin{center}
    \begin{tikzpicture}
        \draw[orange, dashed] (1,-0.5) -- (1,2.5);
        \draw(0,0) rectangle (2,2);
        \draw[->, cyan] (0.5,0.5) -- (1.5,0.5) node[below, midway] {$\small{\vv{B}}$};
        \draw[->, cyan] (0.5,1) -- (1.5,1);
        \draw[->, cyan] (0.5,1.5) -- (1.5,1.5);
        \draw[->] (0,1) -- (0,1.1) node[left] {$i$};
        \draw[->] (2,1) -- (2,0.9) node[right] {$i$};
        \draw[->] (0.9,2) -- (1,2) node[above] {$i$};
        \draw[->] (1.1,0) -- (1,0) node[below] {$i$};
        \node[scale=1.2] at(1,1) {$\bigotimes$};
        \node at(1,1.2)[above] {$\widehat{n}$};
        \draw[->] (1,-0.6) -- (1,-1) node[right] {$\vv{M}$};
    \end{tikzpicture}
    \begin{tikzpicture}
        \draw[->, cyan] (1, 1.8) -- (3,1.8);
        \draw[->, cyan] (1,1.6) ..controls(1.8,1) and(2,1).. (3,1.6) node[below, midway] {$\small{\vv{B}}$};
        \draw[->, cyan] (1,2) ..controls(1.8,2.6) and(2,2.6).. (3,2);
        \draw(0,0) -- (4,0) -- (4,2) -- (3,2) -- (3,1.5) -- (3.5,1.5) -- (3.5,0.5) -- (0.5,0.5) -- (0.5,1.5) -- (1,1.5) -- (1,2) -- (0,2) -- (0,0);
        \node at(0.8,1.8) {$N$};
        \node at(3.2,1.8) {$S$};
        \filldraw (2,1.8) circle (2pt);
        \draw[very thick] (1,1.5) arc (180:0:1);
        \draw[dotted, very thick] (1.16,1.8) arc (160:20:0.9);
        \begin{scope}[rotate = 30, shift={(2.63,0.55)}]
            \draw(-0.5,0) -- (0.5,0);
            \draw(0,0) -- (0,1);
        \end{scope}
    \end{tikzpicture}
\end{center}
Per misurare questa corrente si prende una spira di filo conduttore, che è sospesa tra i poli di un magnete, e quando la corrente attraversa la spira, questa subisce una forza magnetica che tende a farla ruotare, contro bilanciata da una forza elastica che tende a riportarla alla posizione di equilibrio.\\
Questo magnete crea un campo elettromagnetico il cui effetto può essere descritto da un vettore chiamato momento dipolo magnetico: $\vv{m} = i S \widehat{n}$, dove $i$ è la corrente che attraversa la spira, $S$ è l'area (superficie) della spira del galvanometro, e $\widehat{n}$ è un versore perpendicolare alla spira.\\
Questo creerà un momento magnetico che è dato dal prodotto vettoriale tra il momento dipolo magnetico e il campo magnetico creato dal magnete ($N$ e $S$ in figura), e quindi:
\begin{gather*}
        \vv{M} = \vv{m} \land \vv{B}\\
\end{gather*}
Nella fig. a destra il versore $\widehat{n}$ è entrante, quindi il momento magnetico è verso il basso, in arancione c'è lasse di rotazione della spira.\\
La spira ruota dunque in maniera proporzionale alla corrente che la attraversa, consideriamo ora la forza elastica/meccanica che tende a riportare la spira alla posizione di equilibrio:
\begin{gather*}
        \vv{M}_M = -k \varphi \widehat{k}
\end{gather*}
$k$ è la costante elastica del materiale della molla (che cerca di far tornare la spira indietro), $\varphi$ è l'angolo di rotazione della spira, e $\widehat{k}$ è un versore parallelo al momento magnetico ma verso opposto.\\
Otteniamo così un modo per descrivere la forza totale che agisce sulla spira, che è data dalla somma del momento magnetico e del momento meccanico, possiamo dunque 
\begin{gather*}
    \vv{M} = - \vv{M}_M \qquad \Rightarrow \qquad i S \widehat{n} \land \vv{B} =- (- k \varphi \widehat{k})\\
\end{gather*}
Passando ai moduli otteniamo:
\begin{gather*}
    i S B = k \varphi \qquad\Rightarrow\qquad k_r = \frac{k}{S B}\\
    i = k_r \varphi
\end{gather*}
$k_r$ se calcolato è $10^{-6} \frac{A}{rad}$ che è $\lesssim  1000 n A/div$ ed è la costante di proporzionalità del galvanometro, $\varphi$ è sempre l'angolo di rotazione della spira del galvanometro.\\
\par
\begin{wrapfigure}{r}{0.4\textwidth}
    \begin{tikzpicture}
        \draw(0,0) to[R, l=$\rho_G$] (0,1.5) to[short, f=$i$] (0,2) to[rmeter, l = $G$] (0,3);
    \end{tikzpicture}
    \caption{un galvanometro in un circuito}
\end{wrapfigure}
Nei galvanometri di alta precisione (100 nano Ampere che corrispondono a una divisione della scala graduata di riferimento) non c'è un ago fisico (troppo pesante e soggetto ad attriti). Al suo posto si usa un sistema di leva ottica che consiste in uno specchietto montato sull'asse della spira e un raggio luminoso che rimbalza su una scala graduata a grande distanza, anche per rotazioni molto piccole dello specchietto si ha un riscontro sulla scala. 
In generale sono strumenti che hanno una resistenza molto bassa e quindi permettono di misurare correnti molto piccole, ai giorni d'oggi si arriva a misurare correnti di $10^{-6} A$.\\
La curva di taratura di un galvanometro è una retta con coefficiente $\frac{1}{k_r}$ e si ha un'incertezza su di essa sia sul valore misurato che sulla lettura, possiamo rappresentarla così:
\begin{center}
    \begin{tikzpicture}
        \draw[->](0, 0) -- (3, 0) node[at end, below] {$i$};
        \draw[->](0, 0) -- (0, 3) node[at end, left] {$\phi$};
        \draw(0, 0) -- (3, 3) node[at end, right] {$\frac{1}{k_R}$};
        \draw[|-|](0.5, 0.7) -- (0.5, 0.3);
        \draw[|-|](0.3, 0.5) -- (0.7, 0.5);
        \draw[|-|](1.5, 1.7) -- (1.5, 1.3);
        \draw[|-|](1.3, 1.5) -- (1.7, 1.5);
        \draw[|-|](2.5, 2.7) -- (2.5, 2.3);
        \draw[|-|](2.3, 2.5) -- (2.7, 2.5);
    \end{tikzpicture}
\end{center}
\subsection{Confronto con una resistenza standard}
Il galvanometro ha una resistenza interna molto bassa, se confrontiamo questi due circuiti:
\begin{center}
    \begin{tikzpicture}
            \draw(0,1) to [battery1, invert, l=$\mathcal{E}\quad \pm$,] (0,3);
            \draw (0,3) to[R, l=$\rho$] (2,3);
            \draw(2,3) [short, f=$i$] to (2,1);
            \draw(2,1) to [short] (0,1);
        \end{tikzpicture}
        \begin{tikzpicture}
            \draw(0,0) to [battery1, invert, l=$\mathcal{E}\quad \pm$,] (0,2);
            \draw (0,2) to[R, l=$\rho$] (2,2);
            \draw(2,2) to [R, l=$\rho_G$] (2,0);
            \draw(2,0) to [rmeter, l=$G$] (0,0);
    \end{tikzpicture}
\end{center}
E ne calcoliamo la corrente otteniamo:
\begin{gather*}
    i = \frac{\mathcal{E}}{\rho}\\
    i_m = \frac{\mathcal{E}}{\rho+ \rho_G}
\end{gather*}
Se abbiamo che $\rho_G \ll \rho$ allora $i_m \approx i$, ma se abbiamo correnti molto piccole, il loro rapporto con le resistenze conta di più, e quindi $i_m$ è più piccolo di $i$.\\
Di fatto un galvanometro ideale ha resistenza nulla, ma in uno reale dobbiamo trovare un modo per far si che anche con una resistenza interna si possa misurare quasi qualsiasi intensità di corrente.\\
Per ovviare a questo problema si usa quello che si chiama uno Shunt, che è una resistenza 9 volte più piccola di quella del galvanometro.
\begin{center}
    \begin{tikzpicture}
        \draw(-1,1) to[short, f=$i$] (0,1);
        \draw(0,0) to[short] (0,2);
        \draw (0,2) to[R, l=$\rho$, f=$i_G$] (2,2) to[rmeter, l=$G$] (4,2);
        \draw(4,2) to [short] (4,0);
        \draw(4,1) to [short, f=$i$] (5,1);
        \draw(0,0) to [R, l=$\rho_G/G$, f=$i_s$] (4,0);
        \node at(2,-0.2)[below] {Shunt};
    \end{tikzpicture}
\end{center}
Di fatto siccome la corrente passa dove c'è meno resistenza, se metto una resistenza molto piccola in parallelo al galvanometro, la maggior parte della corrente passerà lì, e solo una piccola parte passerà attraverso il galvanometro, che è quello che voglio misurare.\\
Questo ovvia anche al problema di correnti troppo grandi che potrebbero danneggiare il galvanometro.\\
A questo punto abbiamo che la corrente totale sarà $i = i_G + i_s$ ovvero quella del galvanometro più quella dello shunt, sappiamo che il rapporto tra di loro è:
\begin{gather*}
    \frac{i_G}{i_s} = \frac{\rho_G/G}{\rho} = \frac{1}{9}\\
\end{gather*}
Quindi:
\begin{gather*}
    i_s + i_G = i = 10 i_G
\end{gather*}
Da quest'ultima ottengo:
\begin{gather*}
    \boxed{i_G = \frac{1}{10} i}
\end{gather*}
Quindi perdo una sensibilità di un fattore 10, ma riesco a misurare correnti più grandi, e quindi a misurare f.e.m. più grandi.\\
Di fatto lo shunt fa due cose:
\begin{itemize}
    \item Se la corrente è troppo grande, la maggior parte di essa passa attraverso lo shunt, e quindi non danneggia il galvanometro.
    \item Se la resistenza del galvanometro è troppo grande, lo shunt la riduce, e quindi permette di misurare correnti più grandi abbassando la resistenza totale.
\end{itemize}
Se ora calcolo la resistenza equivalente del circuito galvanometro + shunt ottengo:
\begin{gather*}
    R_t = (\frac{1}{\rho G/G} + \frac{1}{\rho})^{-1} = \frac{\cancel{\rho_G} \ \rho_G / G}{\cancel{\rho_G} + \cancel{\rho_G} / G} = \frac{\rho_G / \cancel{G}}{G + \frac{1}{\cancel{G}}} = \frac{\rho_G}{10}
\end{gather*}
Questo fa si che qualsiasi errore sistematico introdotto dal galvanometro venga ridotto di un fattore 10, e quindi si riesce a misurare con più precisione, di fatto:
\begin{gather*}
    i_m = \frac{\E}{\rho + \frac{\rho}{10}} \simeq i
\end{gather*}
\subsection{specifica di linearità di un divisore di tensione}
Torniamo a parlare di divisori di tensione. In particolare ci incentreremo su come funzionano, ricordando generalmente cosa fa un divisore (o partitore) di tensione, data una certa tensione in ingresso $V_{in}$ 
permette di averne una inferiore in uscita $V_{out}$, e può essere regolata con dei cursori.\\
In primo luogo teniamo a mente il concetto dello shunt, che se lo shunt "ruba" una parte della corrente al circuito di fatto il divisore di tensione fa lo stesso con la tensione.\\
Per prima cosa introdurremo il potenziometro che è un dispositivo dotato di una resistenza variabile, e quindi permette di regolare la tensione in uscita secondo la seguente relazione:
\begin{gather*}
    \frac{l \rho}{S} = R
\end{gather*}
Dove $l$ è la lunghezza del potenziometro, $\rho$ è la resistività del materiale di cui è fatto e $S$ è la sezione del potenziometro che attraversa la corrente ad ogni unità di tempo.\\
Il potenziometro è rappresentato sotto, come si vede c'è un cottato mobile su una ruota, in base a quanto si ruota la corrente che scorre è di una lunghezza minore di fatto $R$ è in funzione di $\varphi$.
\begin{center}
    \begin{tikzpicture}
        \draw[very thick](0,1) circle (1);
        \filldraw[page] (-0.5,-0.5) rectangle (0.5,0.5);
        \draw(-2,0.12) to[short] (-0.5,0.12);
        \draw(0.5,0.12) to[short] (2,0.12);
        \draw (-0.5,0.12) to[short] (0,1);
        \filldraw(0,1) circle (2pt);
        \filldraw(-0.5,0.12) circle (2pt) node [below] {$A$};
        \draw[->] (0,1) -- (-0.75,1.65) node[midway, above]{$r$};
        \filldraw(0.5,0.12) circle (2pt) node [below] {$B$};
        \draw (0,1) ++ (140:0.2) arc (140:235:0.2);
        \node at(-0.1,1) [left] {$\varphi$};
    \end{tikzpicture}
\end{center}
La parte spessa rappresenta la resistenza del potenziometro, e il cottato mobile è rappresentato dal punto in alto. Per il teorema dell'asciuga capelli la corrente passa sempre dove c'è meno resistenza, quindi passerà prima dai contatti mobili e poi la strada che gli rimane la percorrerà sulla resistenza.\\
Ora se vogliamo ridurre ad esempio la tensione da $1000 V$ a $1 V$ in modo variabile, servirebbero mille resistenze in serie che si collegano e scollegano o un potenziometro con un filo lunghissimo e metterlo esattamente su di $\frac{1}{1000}$ della sua lunghezza, non molto pratico.\\
Adesso incominciamo anche a occuparci della specifica di linearità di un divisore che è definita come:
\begin{gather*}
    \text{err. linearità} = x \% \cdot \text{ tensioni in ingresso } V
\end{gather*}
Cioò che definisce l'errore è $x$ ma l'errore è proporzionale alla tensione in ingresso, se ad esempio abbiamo $100V$ in ingresso e $x= 0.1\% = \frac{\frac{1}{10}}{100} = \frac{1}{1000}$ allora l'errore sarà di $0.1V$.
Per il potenziometro per un divisore di tensione il rapporto tra $V_{out}$ e $V_{in}$ è in generale $R_{atio} < 1$ (anche perchè il divisore è passivo e non ha batterie interne) ovvero che $V_\text{out} = R_{atio} V_\text{in}$.\\
Se ora vogliamo calcolare l'errore relativo, ovvero $\frac{\Delta V}{V}$, otteniamo:
\begin{gather*}
    \frac{\Delta V}{V} = \frac{\overset{x \%}{\text{linearità }}}{V} = \frac{lin \%}{R_{atio}}
\end{gather*}
Un esempio di un divisre con $20 ppm$ (parti per milione), $R_{atio}=0.5$:
\begin{gather*}
    \frac{\Delta V}{V} = \frac{20 ppm}{0.5} = 40 ppm = 4 \times 10^{-5}
\end{gather*}
\subsection{Divisore di tensione di kelvin-varley}
Il circuito di Kelvin Varley si può schematizzare brevemente su un circuito come illustrato quì sotto, prende in input due rami che hanno una certa differenza di potenziale e alla fine restituisce una tensione in uscita che è una frazione di quella in ingresso.
\begin{center}
    \begin{tikzpicture}
        \draw(0,0) rectangle (2,4);
        \draw(-1,3) -- (0.5,3) node[right] {$H$};
        \draw(-1,1) -- (0.5,1) node[right] {$L$};
        \filldraw(0.5,3) circle(2pt);
        \filldraw(0.5,1) circle(2pt);
        \draw(1.5,3) -- (3,3) node[at start, left] {$h$};
        \draw(1.5,1) -- (3,1) node[at start, left] {$l$};
        \filldraw(1.5,3) circle(2pt);
        \filldraw(1.5,1) circle(2pt);
        \node at(1,2) {$\boxed{r}$};
        \node at(-1,2) {$V$};
        \node at(3,2) {$v$};
    \end{tikzpicture}
\end{center}
Sviluppiamo tutto il circuito per analizzarlo con cura:
\begin{center}
    \begin{tikzpicture}
        \filldraw(-1,0) circle (2pt) node [below] {$L$};
        \filldraw(-1,10.6) circle (2pt) node [above] {$H$};
        \filldraw(1,0) circle (2pt) node [below] {$0$};
        \draw[thick, cyan] (3,4.6) -- (3,2.6) -- (5,2.6);
        \draw[thick, cyan] (7.5,4.6) -- (7.5,2.6) -- (10,2.6);
        \filldraw(7.5,4.6) circle (2pt);
        \filldraw(7.5,8.6) circle (2pt);
        \draw(11,2.6) -- (14,2.6);
        \filldraw(14,2.6) circle (2pt) node [below] {$l$};
        \filldraw(1,2) circle (2pt) node [left] {$1$};
        \filldraw(1,3.96) circle (2pt) node [left] {$2$};
        \filldraw(1,4.6) circle (2pt) node [left] {$8$};
        \filldraw(1,6.6) circle (2pt) node [left] {$9$};
        \filldraw(1,8.6) circle (2pt) node [left] {$10$};
        \draw(-1,0) to[short] (1,0);
        \draw(1,0) to[R] (1,2) to[R] (1,4);
        \node at(1,4) [above]{$\vdots$};
        \draw(1,4.6) to[R] (1,6.6) to[R] (1,8.6) to[R, l=$11 \times 1 K \Omega$] (1,10.6) to[short] (-1,10.6);
        \filldraw(3,8.6) circle (2pt);
        \filldraw(3,4.6) circle (2pt);
        \draw(2.5,4.1) rectangle (3.5,9.1);
        \draw[->] (4,8.6) to (1.1,8.6)  node[above right] {$i+2$};
        \draw[->] (4,4.6) to (1.1,4.6)  node[above right] {$i$};
        \draw(4,8.6) to [short] (4,10.6) to[short] (6,10.6) to[R, l=$11 \times 200 \Omega$] (6,8.6) to[R] (6,6.6) to[R] (6,4.6);
        \draw(4,4.6) to [short] (4,2.6) to[short] (6,2.6) to[R] (6,4) node[above] {$\vdots$};
        \filldraw(5,10.6) circle (2pt) node[above] {$C$};
        \filldraw(5,2.6) circle (2pt) node[below] {$D$};
        \draw[->] (9,4.6) to (6.1,4.6)  node[above right] {$j$};
        \draw[->] (9,8.6)  to (6.1,8.6)node[above right] {$j+2$};
        \draw (7,4.1) rectangle (8,9.1);
        \draw (9, 4.6) to[short] (9,2.6) to[short] (11,2.6);
        \draw (9, 8.6) to[short] (9,10.6) to[short] (11,10.6);
        \filldraw(10,10.6) circle (2pt) node[above] {$E$};
        \filldraw(10,2.6) circle (2pt) node[below] {$F$};
        \draw (11,10.6) node[left]{$10 \times 40 \Omega$} to[R] (11, 8.6) to[R] (11,6.6) to[R] (11,4.6) to[R] (11,2.6);
        \draw[->] (14,6.6) to (11.1,6.6) node[above right] {$k$};
        \filldraw (14,6.6) circle (2pt) node[above]{$h$};
        \path[fill=page] (10.5,4.6) rectangle (11.5,4.2);
        \node at(11,4.5){$\vdots$};
    \end{tikzpicture}
\end{center}
Concentriamoci su un punto alla volta, partendo da sinistra abbiamo 11 resitenze da $1 K \Omega$ l'una in serie tra $H$ e $L$, 
un cursore mobile prende sempre due resistenze tra questo ramo e li mette in parallelo con altre 11 resistenze da $200 \Omega$, quindi un quinto 
della decade precedente, in serie tra $C$ e $D$ per una resistenza totale di $2 k \Omega$, per il teorema dell'asciuga capelli alla corrente conviene dividersi perchè se calcoliamo la resistenza equivalente del ramo $CD$ otteniamo:
\begin{gather*}
    R_{CD} = \frac{2 K\Omega \cdot 2 K\Omega}{2 K\Omega + 2 K\Omega} = 1 K \Omega
\end{gather*}
Quindi il collegamento in parallelo complessivo fa si che ci siano $9$ resistenze da $1 K \Omega$ in serie e una resistenza da $1 K \Omega$ che deriva dal parallelo con l'altra decade (ecco anche perchè si sono messe 11 resistenze e non 10).\\
Si ha poi un'altro ramo con 10 resistenze da $40 \Omega$ in serie che quindi corrispondono a una resistenza da $400 \Omega$, calcoliamone l'equivalente di nuovo:
\begin{gather*}
    R_{EF} = \frac{400 \Omega \cdot 400 \Omega}{400 \Omega + 400 \Omega} = 200 \Omega
\end{gather*}
Quindi ancora una volta alla corrente conviene dividersi, e quindi il secondo ramo ha 11 resistenze ma 2 valgono come una per via del parallelo, infondo al divisore si ha un'ultimo cursore, singolo questa volta, in questo modo 
con i ponti tra un cursore e la base della decade successiva si riesce a ottenere una resistenza variabile che va da $0$ a $1 K \Omega$ con una precisione di $1 \Omega$. Di fatto con i tre cursori possiamo controllare decimi centesimi e millesimi di una resistenza da $1000 \Omega$\\
Vediamo come funziona dal punto di vista della tensione, abbiamo tre differenze di potenziali $V_H - V_L$, $V_C - V_D$ e $V_E - V_F$, e le correnti dei tre cursori $i$, $j$ e $k$.\\
\begin{gather*}
    V_h - V_l = \overbrace{V_h - V_F}^3 + \overbrace{V_F - V_D}^2 + \overbrace{V_D - V_l}^1
\end{gather*}
Procediamo a calcolare la di potenziale della prima decade $V_D - V_l$:
\begin{gather*}
    V_D - V_l = \frac{i \ \cancel{k\Omega}}{ 10 \ \cancel{ k \Omega}} (V_H - V_L) = \frac{i}{10} (V_H - V_L)
\end{gather*}
Ora passiamo alla seconda $V_F - V_D$:
\begin{gather*}
    V_F - V_D = \frac{j}{10}(\frac{1}{10} (V_H - V_L)) = \frac{j}{100} (V_H - V_L)
\end{gather*}
Infine la terza $V_h - V_F$:
\begin{gather*}
    V_h - V_F = \frac{k}{10}(\frac{1}{100} (V_H - V_L)) = \frac{k}{1000} (V_H - V_L)
\end{gather*}
Abbiamo ottenuto che l'equazione iniziale può essere riscritta come:
\begin{gather*}
    1: \quad \frac{i}{10} (V_H - V_L)\\
    2: \quad \frac{j}{100} (V_H - V_L)\\
    3: \quad \frac{k}{1000} (V_H - V_L)
\end{gather*}
Ricapitolando posso riscrivere la tensione in uscita come:
\begin{gather*}
    V_{out} = (\frac{i}{10} + \frac{j}{100} + \frac{k}{1000}) \overbrace{(V_H - V_L)}^{V_{tot}}
\end{gather*}
Ora se prendiamo il cursore $i=5$ $j=3$ e $k=2$ otteniamo una tensione in uscita di $0.532 V$ se la tensione in ingresso è di $1000 V$.\\
\begin{gather*}
    V_{out} = (0.5 + 0.03 + 0.002) \cdot 1000 V = 532 V
\end{gather*}

\newpage
\section{16/3/26}
\subsection{Esperienza numero uno}
fatta per bene su moodle bam bam bam.\\
\subsection{Amplificatore}
È un componenete \textbf{attivo} non come tutti quelli visti fino ad ora


\newpage
\section{18/3/26}
\subsection{Funzionamento dei generatori di corrente}
È un oggetto che cerca di mantenere una certa differenza di potenziale ai suoi capi e si rappresenta inq uesto modo:
\begin{center}
    \begin{tikzpicture}
        \draw(0,0) to[battery1, invert, l=$\E \quad \pm$] (0,2);
    \end{tikzpicture}
\end{center}
Introduciamo ora un fotodiodo che è a sua volta un tipo di generatore di corrente, e si rappresenta inq eusto modo, lo consideriamo come un generatore di corrente ideale
anche questo avrà una sua resistenza interna e idealmente sarebbe schematizzato inq uesto modo:
\begin{center}
    \begin{tikzpicture}
        \draw[->](0,0) -- (3,0) node[below]{$i$};
        \draw[->](0,0) -- (0,3) node[left] {$V$};
        \draw[dashed](2,0) -- (2,2) node[midway, right]{Ideale};
        \draw(2,0) ..controls(1.5,1) and(1.5,1.8).. (1.5,2) node[midway, left]{Reale};
        \draw(0,2) -- (2,2);
        \draw[<->](1.6,1.4) -- (2,1.4) node[midway, above]{$\Delta$};
    \end{tikzpicture}
    \begin{tikzpicture}
        \draw(0,-1) to[led, f=$i$] (0,3);
        \draw(0,0) to[short] (1,0) to[R, l=$R$] (1,2) to[short] (0,2);
        \filldraw (0,2) circle (2pt);
        \filldraw (0,0) circle (2pt);
    \end{tikzpicture}
\end{center}
\begin{center}
    
\end{center}
\begin{gather*}
    i' = I - \frac{\Delta V}{R} \Rightarrow I = i' + \frac{\Delta V}{R}
\end{gather*}

\subsection{Leggi di Norton}
Prendiamo il seguente circuito:
\begin{center}
    \begin{tikzpicture}
        \begin{scope}
            \draw(0,0) to[battery1, invert, l=$\E \quad \pm$, f=$i$] (0,2) to [R, l=$\rho$] (2,2)node[above] {$A$};
            \draw(0,0) to[short] (2,0) node[below]{$B$};
            \filldraw(2,0) circle (2pt);
            \filldraw(2,2) circle (2pt);
        \end{scope}
        \node[scale = 1.5] at(3,1) {$\equiv$};
        \begin{scope}[xshift=5cm]
            \draw(0,0) to[led, l=$\frac{\E}{R}$] (0,2) to[short, f = $i'$] (4,2);
            \draw(2,2) to[short] (2,0) to[R, l=$R$, f=$i''$] (0,0);
            \draw(2,0) to[short] (4,0);
            \filldraw(2,2) circle (2pt);
            \filldraw(4,2) circle (2pt) node[above] {$A$};
            \filldraw(4,0) circle (2pt) node[below] {$B$};
            \filldraw(2,0) circle (2pt);
        \end{scope}
    \end{tikzpicture}
\end{center}
Dimostriamo questa equivalenza.
\begin{gather*}
    \begin{cases}
        V_A - V_B = \frac{\E}{R} \\
    \end{cases}
\end{gather*}
\subsection{Amplificatori operazionali}
\begin{center}
    \begin{circuitikz}
    % Amplificatore Operazionale
    \draw (0,0) node[op amp, scale=1.5
    
    
    
    
    
    
    
    
    
    
    
    
    
    
    ] (opamp) {};
    
    % Ramo di ingresso (Generatore vs e Resistenza Rs)
    \draw (-3, 0.5) to[R, l=$R_s$, i=$ $] (opamp.-);
    \draw (-3, 0.5) -- (-4, 0.5) to[sV, l_=$v_s$] (-4, -1.5) node[ground]{};
    
    % Messa a terra dell'ingresso non invertente
    \draw (opamp.+) -- ++(0,-0.5) node[ground]{};
    
    % Resistenza interna Ri (disegnata tra i due ingressi)
    \draw (opamp.-) to[R, l_=$R_i$] (opamp.+);
    
    % Ramo di retroazione Rf
    \draw (opamp.-) -- ++(0,1.5) coordinate(top) 
          to[R, l=$R_f$, i>=$ $] (top -| opamp.out) -- (opamp.out);
    
    % Nodo di uscita vo
    \draw (opamp.out) to[short, -o] ++(1,0) node[above]{$v_o$};
\end{circuitikz}
\begin{circuitikz}
    % Linea di massa inferiore comune
    \draw (0,0) -- (9,0);
    
    % Generatore di ingresso vs
    \draw (0,0) to[sV, l=$v_s$] (0,3) -- (1,3);
    
    % Resistenza Rs e corrente is
    \draw (1,3) to[R, l=$R_s$, i>=$i_s$] (4,3) coordinate (vin);
    
    % Resistenza di ingresso Ri e tensione v-
    \draw (vin) to[R, l=$R_i$, v=$v_-$] (4,0);
    
    % Resistenza di retroazione Rf e corrente if
    \draw (vin) to[R, l=$R_f$, i>=$i_f$] (7,3) coordinate (vout);
    
    % Ramo di uscita: Resistenza Ro e Generatore dipendente
    \draw (vout) to[R, l=$R_o$] (7,1.5) to[sV, l={$-A v_-$}] (7,0);
    
    % Terminali e nodo di uscita vo
    \draw (vout) to[short, -o] (8.5,3) node[above]{$v_o$};
    \draw (7,0) -- (8.5,0) to[short, -o] (8.5,0);
\end{circuitikz}
\end{center}

\begin{center}
    \begin{circuitikz}
    % Amplificatore Operazionale
    \draw (0,0) node[op amp] (opamp) {};
    
    % Messa a terra dell'ingresso non invertente
    \draw (opamp.+) -- ++(0,-0.5) node[ground]{};
    
    % Nodo comune di ingresso (invertente)
    \coordinate (node) at (-2,0.5);
    \draw (opamp.-) -- (node);
    
    % Tre rami di ingresso
    % Ramo 1: v1, 1k, i1
    \draw (node) |- (-4, 2.5) to[R, l=$1k\Omega$, i=$i_1$] (-6, 2.5)
          to[sV, l=$v_1$, invert] (-6, 1.5) node[ground]{};
          
    % Ramo 2: v2, 10k, i2
    \draw (node) |- (-4, 0.5) to[R, l=$10k\Omega$, i=$i_2$] (-6, 0.5)
          to[sV, l=$v_2$, invert] (-6, -0.5) node[ground]{};
          
    % Ramo 3: v3, R3, i3
    \draw (node) |- (-4, -1.5) to[R, l=$R_3$, i=$i_3$] (-6, -1.5)
          to[sV, l=$v_3$, invert] (-6, -2.5) node[ground]{};
          
    % Ramo di retroazione
    \draw (node) -- ++(0, 1.5) coordinate(top)
          to[R, l=$R_f$, i=$i_f$] (top -| opamp.out) -- (opamp.out);
          
    % Nodo e terminale di uscita vo
    \draw (opamp.out) to[short, -o] ++(1,0) node[above]{$v_o$};
\end{circuitikz}
\end{center}

Integratore
\begin{center}
\begin{circuitikz}
    % Amplificatore Operazionale
    \draw (0,0) node[op amp] (opamp) {};
    
    % Messa a terra dell'ingresso non invertente
    \draw (opamp.+) -- ++(0,-0.5) node[ground]{};
    
    % Ingresso di corrente
    \draw (opamp.-) -- ++(-1.5,0) to[short, i=$I$, -o] ++(-0.5,0);
    
    % Ramo di retroazione con condensatore
    \draw (opamp.-) -- ++(0, 1.5) coordinate(top)
          to[C, l=$C$, i>=$I$] (top -| opamp.out) -- (opamp.out);
          
    % Terminale di uscita
    \draw (opamp.out) to[short, -o] ++(1,0);
\end{circuitikz}
\end{center}


\subsection{Derivatore}
\begin{center}
    \begin{circuitikz}
    % Amplificatore Operazionale
    \draw (0,0) node[op amp] (opamp) {};
    
    % Messa a terra dell'ingresso non invertente
    \draw (opamp.+) -- ++(0,-0.5) node[ground]{};
    
    % Ingresso: Generatore AC e Condensatore C
    \draw (opamp.-) to[C, l_=$C$, i<=$i$] (-3,0.5) 
          to[sV, l_=$v_s$] (-3,-1) node[ground]{};
    
    % Ramo di retroazione con Resistenza R
    \draw (opamp.-) -- ++(0,1.5) coordinate(top)
          to[R, l=$R$, i>=$i$] (top -| opamp.out) -- (opamp.out);
          
    % Terminale e nodo di uscita vo
    \draw (opamp.out) to[short, -o] ++(1,0) node[right]{$v_o$};
    \end{circuitikz}
\end{center}

Tensione di offset
\begin{center}
    \begin{circuitikz}
    % Amplificatore Operazionale
    \draw (0,0) node[op amp] (opamp) {};
    
    % Ingressi N e P all'interno del simbolo (opzionale, per fedeltà al disegno)
    \node at (opamp.-) [right, xshift=0.2cm] {\small N};
    \node at (opamp.+) [right, xshift=0.2cm] {\small P};

    % Ingresso Invertente con Resistenza Rs e Generatore vs
    \draw (opamp.-) to[short, -*] ++(-1,0) coordinate (node_inv)
          to[R, l=$R_s$, i<=$ $] ++(-2.5,0) 
          to[sV, l=$v_s$] ++(0,-1.5) node[ground]{};

    % Generatore di Offset V_offs sull'ingresso non invertente
    \draw (opamp.+) to[battery2, l_=$V_{offs}$] ++(-1.5,0) 
          to[short] ++(0,-1) node[ground]{};

    % Ramo di retroazione con Rf
    \draw (node_inv) -- ++(0,1.5) coordinate(top)
          to[R, l=$R_f$, i>=$ $] (top -| opamp.out) -- (opamp.out);

    % Uscita vo
    \draw (opamp.out) to[short, -o] ++(1,0) node[above]{$v_o$};
    
    \end{circuitikz}
\end{center}

\newpage
\section{15/4/26}
Quest'anno è stata introdotta la nuova esperienza dei cavi di trasmissione.\\
Parte dell'energia si puo perdere per dispersione di energia elettromagnetic.\\\\
In questa lezione introdurremo le grandezze \textbf{alternate} e quindi si incomincierà a parlare di \textbf{corrente alternata}.\\
\textbf{Corrente alternata} è una corrente che cambia direzione periodicamente, una grandezza alternata generica la possiamo descrivere come funione del tempo:
\begin{gather*}
    f(t + T) = f(t) \quad \text{ con } T \text{ periodo}
\end{gather*}
vale che la media di questa funzione su un periodo è 0, ovvero:
\begin{gather*}
    \frac{1}{T} \int_{t}^{t+T} f(t') \ dt' = 0
\end{gather*}
Si incomincierà a parlare di funzioni \textbf{armoniche}, e cominqieremo a parlare di funzioni periodiche che possono descrivere la tensione in questo modo (in funz, del tempo)
\begin{gather*}
    V(t) = V_0 \cos(\omega t + \phi)
\end{gather*}
Ogni grandezza scopriremo che puo essere espressa come serie di Fourier, e quindi come somma di funzioni armoniche. La quale somma troverà la soluzione che cerchiamo.\\
Analizziamo la funzione qua sopra, possiamo definire:
\begin{itemize}
    \item Un ampiezza $V_0$ che rappresenta il massimo valore che la funzione puo assumere
    \item Una frequenza angolare $\omega$ che rappresenta la velocità con cui la funzione oscilla
    \item Uno sfasamento $\phi$ che rappresenta lo spostamento della funzione rispetto all'origine del tempo
\end{itemize}
Il $\cos$ è ciò che rende questa funzione proprio una funzione armonica.\\
\hfill\\
Analizziamo ora il periodo
\begin{gather*}
    T = \frac{2\pi}{\omega}
\end{gather*}
E il periodo inverso che indicheremo con:
\begin{gather*}
    \nu = \frac{1}{T} = \frac{\omega}{2\pi}
\end{gather*}
Che si misura in $s^{-1}$ quindi in $Hz$ ed è proprio la nostra frequenza.\\
\begin{center}
    \begin{tikzpicture}
        \draw[->] (-3.5,0) -- (3.5,0) node[right] {$t$};
        \draw[->] (-pi,-1.2) -- (-pi,1.2) node[above] {$y$};
        \draw[domain=pi:-pi, samples=100, smooth, cyan] plot (\x, {-sin(\x r)});
        \draw(-pi/2+0.1,1) -- (-pi/2-0.1,1) node[above] {$V_0$};
        \draw(pi/2+0.1,-1) -- (pi/2-0.1,-1) node[below] {$-V_0$};
        \draw[orange] (-pi,0) -- (-pi,1) -- (0,1) -- (0,-1) -- (pi,-1) -- (pi,1);
        \draw[yellow] (-pi,0) -- (-pi/2,1) -- (pi/2,-1) -- (pi,0);
        \draw[orange, dashed] (pi,1) -- (pi+0.6,1);
        \draw[yellow, dashed] (pi,0) -- (pi+0.6,0);
    \end{tikzpicture}
\end{center}
Se si iposta un multimetro in AC e si misura tra due cavi impostandolo su voltmetro in AC quello che si misura non è l'ampiezza ma quello che si chiama valore efficace ed è:
\begin{gather*}
    V_{eff} = \sqrt{\frac{1}{T} \int_{t}^{t+T} V^2(t') \ dt'}
\end{gather*}
Ci sono dei generatori che producono forme di onda di ogni tipo a differenza delle esigenze ed hanno un valore efficace diverso, o meglio la differenza tra il valore efficace e l'ampiezza (il picco della funzione) questo dipende dalla forma d'onda.\\
Analizziamone alcuni, partiamo dalla forma d'onda \textbf{armonica} che ci da:
\begin{gather*}
    \frac{V_0}{\sqrt{2}}
\end{gather*}
Mentre il \textbf{quadrato} (detto stepper) ci da semplicemente:
\begin{gather*}
    V_0
\end{gather*}
\hfill\\
Analizziamo ora un po' di componenti, come si comportano in alternata?\\
\begin{itemize}
    \item \textbf{Resistenze} $R$:
    \item \textbf{Induttanze} $L$:
    \item \textbf{Capacità} $C$:
\end{itemize}
In condizioni \textbf{quasi stazionarie} ovvero l'intero circuito evolve contemporaneamente (cioè quando l'intero circuito può essere considerato istante per istante ed è quindi trascurabile la distanza tra i componenti) si può vedere che posso usare piu o meno le stesse leggi che usavo in continua, ma con qualche accorgimento.\\
\begin{center}
    \begin{tikzpicture}
        \draw(0,0) to[vsourcesin] (0,2) to[R, l=$R$] (2,2) to[C, l=$C$] (2,0) to[L, l=$L$] (0,0);
    \end{tikzpicture}
\end{center}
Sappiamo che deve valere la legge delle maglie di Kirchhoff, e quindi:
\begin{gather*}
    \sum_i f_i = \sum_i i_i R_i
\end{gather*}
Quindi che la tensione fornita dal generatore sia uguale alla somma delle cadute di tensione sui singoli componenti, ovvero:
\begin{gather*}
    V(t) = V_R(t) + V_L(t) + V_C(t)
\end{gather*}

Siamo qui in condizioni quasi statiche e stiamo muovendo queste cariche in modo che oscillino di "poco" avrò che:
\begin{gather*}
    \frac{Q(t)}{C} = V(t)
\end{gather*}
Ovvero la tensione che io sto applicando in un dato momento è proporzionale alla carica.\\
Se io ora derivo ottengo uno sfasamento siccome:
\begin{gather*}
    i(t) = \frac{dQ}{dt} = \frac{d}{dt} (V(t) C)\\
    \begin{cases}
        i_0 = V_0 \omega t\\
        \phi = +\frac{\pi}{2}
    \end{cases}
\end{gather*}
Sfruttando il fatto che sto considerando condizioni quasi stazionarie posso ricondurmi alle leggi che già conosciamo dalla continua:
\begin{gather*}
    V_0 \cos(\omega t) = i(t) R + L \frac{d \ i(t)}{d t} + \frac{Q(t)}{C}\\
    \text{sicccome } i(t) = \frac{dQ}{dt} \implies \\ V_0 \cos(\omega t) = L \frac{d^2 Q}{d t^2} + R \frac{dQ}{dt} + \frac{Q(t)}{C}
\end{gather*}
Dove $Q(t)$ è la carica che si accumula sulla capacità, e $i(t)$ è la corrente che scorre nel circuito.\\
Questa è comunque una differenziale di secondo grado difficile da caloclare, perchè ho due incognite la fase e l'ampiezza, la strategia 
che adotterò sarà quella di considerare $Q$ e $i$ come numeri complessi, e quindi ricondurmi a una differenziale di primo grado, che è molto più semplice da risolvere.\\
%\begin{gather*}
%    \tilde{V} = V_0 (\cos(\cos(\omega t)) + j \sin(\omega t)) = V_0 e^{j \omega t}\\
%    \text{considerando la fase:}\\
%    \tilde{V} = V_0 e^{j (\omega t + \varphi)} = V_0 e^{j \varphi} e^{j \omega t}
%\end{gather*}
%\hfill\\
%I numeri complessi praticamente ci permettono di esprimere con la forma trigonmetrica due variabili coem una, e quindi ricondurci ai casi che sappiamo gestire.\\
%\begin{gather*}
%    V(t) = i(t) R\\
%    V_0 \cos(\omega t) R\\
%    \begin{cases}
%        i_0 = \frac{V_0}{R}\\
%        \varphi = 0
%    \end{cases}
%\end{gather*}
%Fare attensione che la fase la sto considerando 0, ma la sto sempre scrivendo, tenendo in considerazione, questo perchè se avessi una induttanza o una capacità la fase non sarebbe 0, e quindi è importante tenere traccia di questa variabile.\\
Prima di fare questo però analizziamo il circuito nei suoi vari pezzi, guardando singolarmente ogni componente.\\
\subsection*{Resistenza:}
\begin{center}
    \begin{tikzpicture}
        \draw(0,0) to[vsourcesin] (0,2) to[R, l=$R$] (2,2) to[short] (2,0) to[short] (0,0);
    \end{tikzpicture}
\end{center}
Sempre per la quasi staticità dei componenti mi riconduco alla legge di ohm, e quindi:
\begin{gather*}
    V(t) = i(t) R
\end{gather*}
Sappiamo che il generatore oscilla come un'onda cosinusoidale, conquindi pulsazione $\omega$ e ampiezza $V_0$, possiamo quindi scrivere la tensione come:
\begin{gather*}
    V(t) = V_0 \cos(\omega t)
\end{gather*}
Da questa possiamo mettere in relazione corrente e tensione, isolando la corrente otteniamo:
\begin{gather*}
    V_0 \cos(\omega t) = i(t) R\\
    i(t) = \frac{V_0}{R} \cos(\omega t)
\end{gather*}
Abbiamo ricavato che:
\begin{itemize}
    \item \textbf{L'ampiezza}: il valore massimo della corrente è smeplicemente $i_0 = \frac{V_0}{R}$
    \item \textbf{La fase}: L'argomento del coseno per la tensione è ($\omega t$), e per la corrente è sempre ($\omega t$). Non c'è \textbf{nessuno sfasamento}.
\end{itemize}
\begin{gather*}
    \begin{cases}
        i_0 = \frac{V_0}{R}\\
        \varphi = 0
    \end{cases}
\end{gather*}
\subsection*{Capacità}:
\begin{center}
    \begin{tikzpicture}
        \draw(0,0) to[vsourcesin] (0,2) to[C, l=$C$] (2,2) to[short] (2,0) to[short] (0,0);
    \end{tikzpicture}
\end{center}
Il generatore ha la solita legge per la tensione:
\begin{gather*}
    V(t) = V_0 \cos(\omega t)
\end{gather*}
Fisicamente, sappiamo che in un condensatore la quantità di carica Q accumulata sulle armature è direttamente proporzionale alla tensione applicata in quel momento:
\begin{gather*}
    Q(t) = C V(t)
\end{gather*}
Noi però vogliamo trovare la corrente i(t). Per definizione, la corrente non è altro che la rapidità con cui si muovono le cariche, ovvero la derivata nel tempo della carica:
\begin{gather*}
    i(t) = \frac{dQ}{dt}
\end{gather*}
Ora mettiamo tutto insieme. Sostituiamo Q(t) nella formula della corrente e deriviamo:
\begin{gather*}
    i(t) = \frac{d}{dt} (C V(t)) = C \frac{d}{dt} (V_0 \cos(\omega t)) =\\
     - C V_0 \omega \sin(\omega t)\\
     \text{da identità:} \implies C V_0 \omega \cos(\omega t + \frac{\pi}{2})
\end{gather*}
\begin{itemize}
    \item \textbf{L'ampiezza}: $i_0 = C V_0 \omega$, più la frequenza è alta più corrente passa.
    \item \textbf{La fase}: L'argomento del coseno per la tensione è ($\omega t$), e per la corrente è ($\omega t + \frac{\pi}{2}$). Quindi c'è uno sfasamento di $\frac{\pi}{2}$.
\end{itemize}
\begin{gather*}
    \begin{cases}
        i_0 = C V_0 \omega\\
        \varphi = +\frac{\pi}{2}
    \end{cases}
\end{gather*}
\subsection*{Induttanza:}
\begin{center}
    \begin{tikzpicture}
        \draw(0,0) to[vsourcesin] (0,2) to[L, l=$L$] (2,2) to[short] (2,0) to[short] (0,0);
    \end{tikzpicture}
\end{center}


\hfill\\
Se prendiamo un solenoide e ci facciamo passare una corrente alternata, creeremo un campo magnetico $B$:
\begin{center}
    \begin{tikzpicture}[scale = 3, transform shape]
        \draw(0,0) to[L, l=\tiny{$L$}] (2,0);
        \draw[->,ultra thick, cyan](0.8,0.05) -- (1.2,0.05) node[midway, below]{\tiny{$B$}};
    \end{tikzpicture}
    $\qquad$
    \begin{tikzpicture}
        \draw(0,0) circle (1);
        \draw[->](-0.1,1) -- (0.1,1) node[above] {$f$};
        \node at (0,0) {\huge{$\odot$}};
        \node at(0,-0.2) [below] {$\vv{B}$};
    \end{tikzpicture}
\end{center}
Sappiamo che il generatore fornisce sempre la stessa tensione quindi:
\begin{gather*}
    V(t) = V_0 \cos(\omega t)
\end{gather*} 
Cosa succede nell'induttanza? Quando una corrente variabile la attraversa, per la legge di Faraday-Neumann-Lenz si genera una forza elettromotrice indotta che si oppone alla variazione di corrente. Questa è data da:
\begin{gather*}
    {f.e.m.}_\text{indutt.} = \frac{d \varphi(B)}{dt} \qquad \text{con: }\varphi(B) = L I\\
    {f.e.m.}_\text{indutt.} = -L \frac{d_i}{d_t}\\
\end{gather*}
Di fatto un'induttanza è a tutti gli effetti un solenoide\\
\hfill\\
Tornando al circuito, possiamo mettere in relazione la tensione del generatore con la forza elettromotrice indotta, sfruttando l'uguaglianza del generatore, otteniamo:
\begin{gather*}
    V(t) = -L \frac{d_i}{d_t} = 0\\
    V_0 \cos(\omega t) -L \frac{d_i}{d_t} = 0\\
    \frac{d_i}{d_t} = \frac{V_0}{L} \cos(\omega t)
\end{gather*}
Il nostro obbiettivo è sempre trovare $i(t)$, quindi isoliamo la corrente integrando rispetto al tempo:
\begin{gather*}
    \int \frac{d_i}{d_t} dt = \int \frac{V_0}{L} \cos(\omega t) dt\\
    \Rightarrow i(t) = \frac{V_0}{L \omega} \sin(\omega t)
\end{gather*}
\begin{itemize}
    \item \textbf{L'ampiezza}: $i_0 = \frac{V_0}{L \omega}$, di nuovo più la frequenza è alta meno corrente passa.
    \item \textbf{La fase}: L'argomento del coseno per la tensione è ($\omega t$), e per la corrente è ($\omega t - \frac{\pi}{2}$). Quindi c'è uno sfasamento di $-\frac{\pi}{2}$.
\end{itemize}
\begin{gather*}
    \begin{cases}
        i_0 = \frac{V_0}{L \omega}\\
        \varphi = -\frac{\pi}{2}
    \end{cases}
\end{gather*}
Che è la relazione della corrente per il solenoide.\\
\section*{Risoluzione con i numeri complessi}
Riprendiamo ora il circuito con resistenza, induttanza e capacità:
\begin{center}
    \begin{tikzpicture}
        \draw(0,0) to[vsourcesin] (0,2) to[R, l=$R$] (2,2) to[C, l=$C$] (2,0) to[L, l=$L$] (0,0);
    \end{tikzpicture}
\end{center}
Ricordiamo che siamo in condizioni qusi statiche e quindi valgono le leggi viste prima, richiamiamole tutte insieme:
\begin{gather*}
    V(t) - L \frac{d_i}{d_t} - \frac{Q(t)}{C} = i(t) R\\
    i = \frac{dQ}{dt} \Rightarrow V(t) = L \frac{d^2 Q}{d t^2} + \frac{Q(t)}{C} + R \frac{dQ}{dt}
\end{gather*}
Questa è una equazione differenziale di secondo grado, non è particolarmente difficile da risolvere, ma\\
è molto lunga, quà riportiamo direttamente i risultati:
\begin{gather*}
\begin{cases}
    Q_0 = -\frac{V_0}{\omega \sqrt{R^2 + (\omega L - \frac{1}{\omega C})^2}}\\
    \tan \varphi = \frac{R}{\omega L - \frac{1}{\omega C}}
\end{cases}
\end{gather*}
Questo risultato che otteniamo con il metodo "classico" prende circa due pagine di conti.\\
Per evitarci tutto questo, introduciamo il metodo simbolico (i fasori) tramite i numeri complessi e la formula di Eulero:
\begin{gather*}
    V(t) = V_0 \cos(\omega t) \implies \tilde{V} = V_0 (\cos(\omega t) + j \sin(\omega t)) = V_0 e^{j \omega t}
\end{gather*}
Se teniamo conto di uno sfasamento $\varphi$ otteniamo:
\begin{gather*}
    \tilde{V} = V_0 \cos(\omega t + \varphi) \implies \tilde{V} = V_0 e^{j (\omega t + \varphi)} = V_0 e^{j \varphi} e^{j \omega t}
\end{gather*}
Quello che stiamo facendo è considerare una parte immaginaria (che poi butteremo via) per sfruttare la forma trigonometrica 
dei numeri complessi (dalla formula di Eulero), questo ci semplifica notevolmente i conti che per lo più saranno derivate e integrali (banali quando trattiamo esponenziali).\\
Una volta risolta l'equazione differenziale, ci basterà prendere la parte reale del risultato per ottenere la soluzione che cerchiamo.\\
Così facendo generalizziamo la legge di ohm come:
\begin{gather*}
    \tilde{V} = \tilde{I} Z
\end{gather*}
Con $Z$ che è l'impedenza del circuito, e che è data da:
\begin{gather*}
    \begin{cases}
        Z_R = R\\
        Z_C = \frac{1}{j \omega C}\\
        Z_L = j \omega L
    \end{cases}
\end{gather*}
le impedenze si sommano come le resistenze quindi;
\begin{gather*}
    z_s = \sum_i z_i\\
    \frac{1}{z_p} = \sum_i \frac{1}{z_i}
\end{gather*}

%Come si comporta il circuito in funzione della frequenza?\\
%Se vado a frequenza zero vedo che l'impedenza di una capacità\\
%\begin{gather*}
%    \tilde{V} = \tilde{I} (R + \frac{1}{j \omega C} + i\omega L)\\
%\end{gather*}
%Questa è la legge di ohm generalizzata a patto che tu abbia definito le impedenze nel modo che abbiamo fatto vedere prima.\\
%Si considera anche $Q$ un numero complesso.
%leo ha fatto na domanda:\\
%prendiamo il primo esempio
%\begin{center}
%    fig
%\end{center}
%\begin{gather*}
%    \tilde{V} = \tilde{I} R\\
%    \tilde{V} = V_0 e^{j \omega t}\\
%    \tilde{I} = \frac{V_0}{R}
%\end{gather*}
%Questo mi diche che io ho un numero complesso che ha:
%\begin{gather*}
%    |\tilde{I}| = \frac{V_0}{R}\\
%    \varphi = 0
%\end{gather*}
%Ora con l'esempio della capacità:
%\begin{center}
%    fig
%\end{center}
%\begin{gather*}
%    \tilde{V} = \tilde{I} \frac{1}{j \omega C}\\
%    \tilde{I} = j \omega C V_0 \begin{cases}
%        V_0 \omega C = |I|\\
%        \varphi = +\frac{\pi}{2}
%    \end{cases}
%\end{gather*}
%*nuova punchline* "BAM trovato!"\\
\hfill\\
Tornando all'esempio del solenoide:
%\begin{center}
%    \begin{tikzpicture}
%        \draw(0,0) to[vsourcesin] (0,2) to[R, l=$R$] (2,2) to[C, l=$C$] (2,0) to[L, l=$L$] (0,0);
%    \end{tikzpicture}
%\end{center}
\begin{gather*}
    \tilde{V} = \tilde{I} j \omega L\\
    I = \frac{V_0}{j \omega L} = - j\frac{V_0}{\omega L} \begin{cases}
        |I| = \frac{V_0}{\omega L}\\
        \varphi = -\frac{\pi}{2}
    \end{cases}
\end{gather*}
Riprendiamo ora il circuito con resistenza, induttanza e capacità:
\begin{center}
    \begin{tikzpicture}
        \draw(0,0) to[vsourcesin] (0,2) to[R, l=$R$] (2,2) to[C, l=$C$] (2,0) to[L, l=$L$] (0,0);
    \end{tikzpicture}
\end{center}
\begin{gather*}
    \tilde{V} = \tilde{I} (R + \frac{1}{j \omega C} + j\omega L)\\
\end{gather*}
Ricordando che $j$ funge da numero immaginario e quindi $\frac{1}{j} = -j$, la possiamo anche riscrivere come:
\begin{gather*}
    Z = R + \frac{1}{j \omega C} + j\omega L = \\
    R  - j\frac{1}{\omega C} + j\omega L = \\
    R + j (\omega L - \frac{1}{\omega C})
\end{gather*}
\begin{gather*}
    \tilde{I} = \frac{V_0}{R + j (\omega L - \frac{1}{\omega C})} = \frac{V_0(R - j(\omega L - \frac{1}{\omega C}))}{R^2 + (\omega L - \frac{1}{\omega C})^2}\\
    R_e = \frac{V_0 R}{(R^2 + (\omega L - \frac{1}{\omega C})^2)} \qquad I_m = \frac{(\omega L -\frac{1}{\omega C}) V_0}{(R^2 + (\omega L - \frac{1}{\omega C})^2)}\\
\end{gather*}
Ora facciamo il modulo di $\tilde{I}$ (quindi $\sqrt{Re^2 + Im^2}$) e la tangente dell'argomento:
\begin{gather*}
    |I| = \sqrt{\frac{V_0^2 \ R^2}{{(R^2 + (\omega L - \frac{1}{\omega C})^2)^{\cancel{2}}} } + \frac{V_0^2 (\omega L - \frac{1}{\omega C})^2}{{(R^2 + (\omega L - \frac{1}{\omega C})^2)^{\cancel{2}}} }} = \sqrt{\frac{V_0^2 \cancel{(R^2 + (\omega L - \frac{1}{\omega C})^2)}}{(R^2 + (\omega L - \frac{1}{\omega C})^2)^{\cancel{2}}}} = \frac{V_0}{\sqrt{R^2 + (\omega L - \frac{1}{\omega C})^2}}\\
\end{gather*}
Sempre dalla rappresentazione trigonometrica dei numeri complessi possiamo ricavare la fase che è l'angolo tra la parte reale e quella immaginaria, siccome conosciamo già i moduli delle due componenti:
\begin{gather*}
    \tan(\varphi) = \frac{I_m}{R_e}
    \tan(\varphi) = \frac{(\omega L - \frac{1}{\omega C}) \cancel{V_0}}{\cancel{V_0} R} = \frac{\omega L - \frac{1}{\omega C}}{R}
\end{gather*}
"PAPPARDELLA" $= R^2 + (\omega L - \frac{1}{\omega C})^2$\\



\newpage
\section{20/04/26}
Prendiamo un semplice esempio in cui abbiamo un condensatore e una resistenza, con in mezzo un'interruttore.
\begin{center}
    \begin{tikzpicture}
        \draw(0,0) to[short] (2,0) to[R, l=$R$] (2,2) to[short] (0,2) to[C, l=$C$, a_=$-Q_0$] (0,0);
    \end{tikzpicture}
\end{center}
In questo circuito non c'è ne irraggiamento ne nessuna resistenza sul condensatore.
Facciamo il conto di cosa succede in condizioni quasi stazionarie, chiudendo il circuito:
\begin{gather*}
    \frac{Q(t)}{C} = V(t)
\end{gather*}
Applico la seconda legge di kirchhoff e ottengo:
\begin{gather*}
    V(t) = i(t) R
\end{gather*}
Ad ogni variazione di carica corrisponderà una variazione di corrente positiva, attenzione merò al segno meno sulla carica:
\begin{gather*}
    -dQ = I(t) dt\\
    I(t) = - \frac{dQ(t)}{dt}
\end{gather*}
Continuando:
\begin{gather*}
    \frac{Q(t)}{C} = -\frac{dQ(t)}{dt} \ R\\
    \int \frac{dQ(t)}{dt} = \int -\frac{1}{R C} dt\\
    \begin{cases}
        Q(t) = Q_0 e^{-\frac{t}{R C}}\\
        I(t) = \frac{Q_0}{R C} e^{-\frac{t}{R C}}
    \end{cases}
\end{gather*}
Caolcolando ora la f.e.m. indotta:
\begin{gather*}
    \E = \frac{1}{2} \int_V \E_0 t^2 dV = \frac{1}{2} \frac{Q_0^2}{C} \overset{?}{=} \int^{\infty}_{0} i^2(t) R dt
    V = \int \vv{\E} \ \vv{d\ell} = E \cdot d \implies E = \frac{V}{d}\\
    \frac{Q}{C} = V \qquad E = \frac{Q}{C d}
\end{gather*}
potenziale...
\begin{gather*}
    U_t = \int_0^\infty i^2(t) R dt = \int_0^\infty \frac{Q_0^2}{(R C)^2} e^{-\frac{2t}{R C}} R dt \\
    = \frac{Q^2}{R^2 C^2} R \int_0^\infty e^{-\frac{2t}{R C}} dt = \frac{Q^2}{R C^2} \left[ -\frac{RC}{2} \right] \left[ e^{-\frac{2t}{RC}} \right]_0^\infty = \\
    = \frac{Q^2}{\cancel{R} C^{\cancel{2}}} (-\frac{\cancel{R}\cancel{C}}{C}) (0-1) = \frac{1}{2} \frac{Q^2}{C}\\
    \int e^{-\alpha x} dx = -\frac{1}{\alpha} e^{-\alpha x}
    \begin{cases}
        y = \alpha x\\
        dy = \alpha dx\\
    \end{cases}
\end{gather*}
\begin{example}
    esempio pratico\\
    \begin{gather*}
        R = 100 \Omega = 10^2 \Omega\\
        C = 0,1 \mu F = 10^{-7} F\\
        \tau = 10^2 \cdot 10^{-7} s = 10^{-5} s = 10 \mu s \qquad \tau \gg \Delta \\
        \ell = 10 cm \qquad \Delta t = \frac{\ell}{C} = \frac{10 \ 10^{-2} m}{3 \ 10^8 m/s} = \frac{10}{3} \cdot 10^{-10} s = 300 ps
    \end{gather*}
    (domanda di marco fabbri figura di uno schema di un'antenna)
\end{example}
Il bro sta letteralmente facendo un esercizio di fisica due perchè lo "ha chiesto la classe" (lo ha chiesto lui alla classe e nessuno lo ha negato)

\section*{Potenza nei circuiti elettrici}
Prendiamo un percorso che compie una carica e va da un punto $A$ ad un punto $B$.
\begin{gather*}
    d \mathcal{L} = dq \ \vv{E} \cdot \vv{dl} = n dV q \ \vv{E} \cdot \vv{dl}\\
    \vv{dl} = \frac{\vv{v}_d}{\vv{J}} \ dt\\
    d\mathcal{L} = (n \vv{v} \cdot q) \vv{E} dV \ dt\\
    \frac{W}{m^3} = \frac{J}{m^3 s} = \frac{d \mathcal{L}}{dV dt} = \vv{E} \cdot \vv{J}\\
    dL = \int_A^B dq \vv{E} \cdot \vv{dl} = dq (V_A - V_B) = I(V_A - V_B) dt\\
    dq = I dt \qquad W = \frac{dL}{dt} = I \Delta V = \frac{\Delta V^2}{R} = I^2 R
\end{gather*}
Ora sfruttiamo il metodo simbolico che abbiamo introdotto la scorsa volta (che semplifica notevolmente i conti)
\begin{gather*}
\begin{cases}
        V(t) = V_0 \cos(\omega t)\\
        I(t) = I_0 \cos(\omega t + \varphi)
\end{cases} \implies 
\begin{cases}
        \tilde{V} = V_0 e^{j \omega t}\\
        \tilde{I} = I_0 e^{j \omega t + \varphi} = I_0 e^{j \varphi} e^{j \omega t} = \tilde{I}_0 e^{j \omega t}
\end{cases}
\end{gather*}
Quindi ritroviamo la legge di ohm geenralizzata:
\begin{gather*}
    \tilde{V}_0 = \tilde{I}_0 Z
\end{gather*}
E avevamo visto che vale:
\begin{gather*}
    Z_R = R\\
    Z_C = \frac{1}{j \omega C}\\
    Z_L = j \omega L
\end{gather*}
Partiamo da questo fatto:
\begin{gather*}
    Re \{ \tilde{a} \cdot \tilde{b} \}  \neq Re\{\tilde{a}\} \cdot Re\{\tilde{b}\}
\end{gather*}
Quindi sul calcolo delle potenze devo stare attento\\
\begin{center}
    fig
\end{center}
\begin{gather*}
    z = z_0 e^{-j \varphi}
\end{gather*}
\begin{gather*}
    \begin{cases}
        V(t) = V_0 \cos(\omega t)\\
        I(t) = I_0 \cos(\omega t + \varphi)
    \end{cases}
    \implies W(t) = I(t) V(t) = V_0 i_0 (\cos(\omega t) \cos(\omega t + \varphi)) = \\
    = V_0 i_0 \cos(\omega t) \ \left( \ \cos(\omega t) + \cos(\varphi) - \sin(\omega t) \sin(\varphi) \ \right)\\
    W(t) = V_0 i_0 \left( \cos^2(\omega t) \cos(\varphi) - \cos(\omega t)\sin(\omega t) \sin(\varphi) \right)
\end{gather*}
Calcoliamo ora la potenza media su un periodo:
\begin{gather*}
    \left\langle W(t) \right\rangle = \frac{V_0 i_0}{2} \cos(\varphi) = \frac{V_0}{\sqrt{2}} \frac{i_0}{\sqrt{2}} \cos(\varphi) = V_{eff} I_{eff} \cos(\varphi)
\end{gather*}
Quindi potenza dissipata la ho ma solo in alcuni casi:
\begin{gather*}
    \begin{cases}
        \text{se } \varphi = 0 \implies W_{diss} = V_{eff} I_{eff}\\
        \text{se } \varphi = \pm \frac{\pi}{2} \implies W_{diss} = 0
    \end{cases}
\end{gather*}
L'energia passa dal campo elettrico al campo magnetico senza disspiarsi, ovviamente avrò che:
\begin{gather*}
    \left\langle W(t) \right\rangle \neq 0 \qquad \text{ se } \cos(\varphi) \neq 0 \implies Re\{ \tilde{z} \} \neq 0
\end{gather*}
\section*{Filtri}
Ne vedremo di due tipologie passa alto e passa basso.\\
\begin{center}
    Filtro passa basso:\\
    \begin{tikzpicture}
        \draw(0,0) to[vsourcesin, l=$V_1(t)$] (0,2) to[R, l=$R$] (2,2) to[short] (2,0) to[C, l=$C$] (0,0);
        \draw(0,0) to[short] (0,-1);
        \draw(2,0) to[short] (2,-1);
        \node at (1,-1.5) {$V_2(t)$};
    \end{tikzpicture}
\end{center}
Mi interessa studiare il rapporto $\frac{\tilde{V}_2}{\tilde{V}_1}$, che sarà proprio:
\begin{gather*}
    \frac{\tilde{V}_2}{\tilde{V}_1} = \frac{\frac{1}{j \omega C}}{R + \frac{1}{j \omega C}} = \frac{1}{1 + j \omega R C} = \frac{1}{1+j \frac{\omega}{\omega_0}}
\end{gather*}
Analizziamo ora le dimensioni ottenute da questi calcoli:
\begin{gather*}
    R C = [ s ]\\
    \frac{1}{RC} = [\frac{rad}{s}] = \omega_0\\
    \omega = 2 \pi f
    \omega_0 = 2 \pi f_0
\end{gather*}
Ora possiamo analizzare l'ampiezza:
\begin{gather*}
    \tilde{A} (f) = \frac{1}{1+ j \frac{f}{f_0}}\\
    \tilde{A} (f) = \frac{1}{1 + j \frac{f}{f_0}} \cdot \frac{1 - j \frac{f}{f_0}}{1 - j \frac{f}{f_0}} = \frac{1 - j \frac{f}{f_0}}{1 + (\frac{f}{f_0})^2}\\
    \begin{cases}
        Re = \frac{1}{1+ (\frac{f}{f_0})^2} = a\\
        Im = - \frac{\frac{f}{f_0}}{1 + (\frac{f}{f_0})^2} = b
    \end{cases}
\end{gather*}
Quindi ottengo che:
\begin{gather*}
    |\tilde{A}| = \frac{1}{\sqrt{1 + (\frac{f}{f_0})^2}} \begin{cases}
        1 \quad \text{per } f \ll f_0\\
        \frac{1}{\sqrt{2}} \quad \text{per } f = f_0\\
        0 \quad \text{per } f \gg f_0
    \end{cases}
\end{gather*}
Ora utilizziamo i decibel per rappresentare l'ampiezza, quindi $|\tilde{A}|_\text{dB}$, il decibel è:
\begin{gather*}
    dB = 20 \log_{10} \frac{V_2}{V_1}\\
    dB = 10 \log_{10} \frac{P_2}{P_1} \qquad \text{con } P = \frac{V^2}{R}
\end{gather*}
Quindi:
\begin{gather*}
    |\tilde{A}|_\text{dB} = 20 \log_{10} |\tilde{A}| = 20 \log_{10} \frac{1}{\sqrt{1 + (\frac{f}{f_0})^2}} = -10 \log_{10} (1 + (\frac{f}{f_0})^2)\\
    \begin{cases}
        0 \quad \text{per } f \ll f_0\\
        \simeq -3 \quad \text{per } f = f_0\\
        \begin{matrix}
            -20 dB \text{ per decade} \\
            -6 dB \text{ per ottava}
        \end{matrix}\quad \text{per } f \gg f_0
    \end{cases}
\end{gather*}
Facciamo il grafico di questa ampiezza:
\begin{center}
    \begin{tikzpicture}
        \draw[->](0,0) -- (5,0) node[below] {$f$};
        \draw[->](0,0) -- (0,3) node[left] {$|\tilde{A}|_\text{dB}$};
        \draw(1,0) .. controls (1.5,-0.2) and (3,-0.5) .. (4,-2) node[below] {$a$};
    \end{tikzpicture}
\end{center}
Questo ci da la freq. di taglio che non è abbastanza per descrivere ciò che ci interessa.
\begin{gather*}
    \tan(\varphi) = \frac{Im}{Re} = - \frac{f}{f_0}\\
    \varphi = - \arctan(\frac{f}{f_0})
    \begin{cases}
        0 \quad \text{per } f \ll f_0\\
        -\frac{\pi}{4} \quad \text{per } f = f_0\\
        -\frac{\pi}{2} \quad \text{per } f \gg f_0
    \end{cases}
\end{gather*}



\end{document}